\documentclass[11pt,reqno]{amsart}
\usepackage[a4paper,margin=27mm]{geometry}
\usepackage[T1]{fontenc}
\usepackage{lmodern}
\usepackage{amsmath,amssymb,amsthm,mathtools,mathrsfs,booktabs,longtable,array,enumitem,microtype}
\usepackage[colorlinks=true,linkcolor=blue!50!black,citecolor=blue!50!black,urlcolor=blue!50!black]{hyperref}
\usepackage{xcolor}
\usepackage[nameinlink,capitalise,noabbrev]{cleveref}
\hypersetup{pdftitle={Transport energy of the prime-corrected Riemann--von Mangoldt clock},pdfauthor={Rotkiv Aynasac},pdfsubject={Riemann zeta zeros, Euler phases and corrected clocks}}
\numberwithin{equation}{section}
\newtheorem{theorem}{Theorem}[section]
\newtheorem{proposition}[theorem]{Proposition}
\newtheorem{lemma}[theorem]{Lemma}
\newtheorem{corollary}[theorem]{Corollary}
\newtheorem{definition}[theorem]{Definition}
\newtheorem{remark}[theorem]{Remark}
\newcommand{\R}{\mathbb R}
\newcommand{\C}{\mathbb C}
\newcommand{\N}{\mathbb N}
\newcommand{\av}[1]{\left\langle #1\right\rangle_T}
\newcommand{\norm}[1]{\left\|#1\right\|_{2,T}}
\newcommand{\Li}{\operatorname{Li}}
\newcommand{\e}{\mathrm e}
\newcommand{\calI}{\mathcal I_T}
\newcommand{\lamb}{\Lambda}
\newcommand{\eps}{\varepsilon}
\newcommand{\Cloc}{C_{\mathrm{loc}}}
\newcommand{\avgJ}[1]{\left\langle #1\right\rangle_{\mathcal J_T}}
\title[Transport energy of the corrected clock]{Transport energy of the prime-corrected Riemann--von Mangoldt clock and genuinely index-dependent prime cutoffs}
\author{Rotkiv Ayn\"asac}
\date{24 October 2026}
\begin{document}
\begin{abstract}
The midpoint Riemann--von Mangoldt clock is corrected by a finite Euler phase formed from prime powers, without using ordinates of the zeta zeros as construction inputs. We obtain exact multiplicity-preserving identities for the discrepancy at the true ordinates, a uniform second-moment estimate for the residual prime information, and an asymptotic for the transport energy of the one-step clock. For every fixed $0<c\le1/2$, the squared prime-phase transport averaged over ordinates in $(T,2T]$ equals $\frac1{2\pi^2}\log\log\log T+O_c(\sqrt{\log\log\log T\,\log\log\log\log T})$. The corresponding corrected ordinate error is $[2\log\log\log T+O_c(\sqrt{H_1H_2})]/\log^2T$. We establish uniform moving-cutoff consequences, a harmonic cancellation estimate for the bare-clock response, and an arithmetic characteristic profile for the explicitly computable second-clock response over arbitrary prime subsets. The proof is based on corrected-point multiplicity, carefully bounded prime-power truncations, a two-scale aggregated Taylor formula, and deterministic bounds for every nonzero rational frequency. We retain the endpoint $c=1/2$, full zero multiplicities and boundary corrections. Further unconditional theorems prove operator-norm almost orthogonality for prime-power phases up to $T^b$, for each fixed $b<1$, exact S2 orthogonal decomposition, and stability of the orthogonal remainder for genuinely index-dependent one-step cutoffs constrained only by a thin harmonic prime band. These estimates transfer the fixed-power energy law to arbitrarily oscillating indexwise multipliers. We give explicit counterexamples distinguishing scalar energy from operator-norm orthogonality, positive error variance from failed localization, and fixed-exponent limits from the endpoint $b\uparrow1$; a boundary stationary-phase calculation also rules out a frequently used uniform full-Gaussian formula.
\end{abstract}
\maketitle
\tableofcontents
\section{Introduction}
The smooth term in the Riemann--von Mangoldt zero counting formula supplies an explicit increasing clock whose $k$th point is prescribed by the half-integer equation. The true ordinates fluctuate around this clock. Prime sums approximate the fluctuating part of the zero count, but their values at the true ordinates and at clock points are different. The difference is not a negligible formal substitution: moving the argument of a trigonometric polynomial whose largest frequency grows like $\log T$ produces a nonlinear phase response.

The present analysis separates three distinct operations. First, information about primes is evaluated at a true ordinate. Second, the same finite prime-power phase is sampled at a deterministic clock point. Third, a correction formed from that clock value is compared to the ordinate. These operations have different quantitative errors. Their exact relation is given below without assuming the Riemann hypothesis and without requiring monotonicity of the prime-corrected level function.

We use a midpoint clock and all positive ordinates of zeta zeros, counted with multiplicity. In particular $\rho_k=\beta_k+i\gamma_k$ and $0<\beta_k<1$ throughout; no condition $\beta_k=1/2$ enters any of the definitions. Our algebraic lemmas apply even when consecutive ordinates coincide.

A key technical distinction concerns what can be concluded from a second moment. An energy estimate for corrected ordinates gives control of typical displacements; it does not yield a pointwise separation theorem. Likewise the prime-phase characteristic function of an independent model supplies the diagonal of certain polynomial expansions, not automatically the complete deterministic response at the clock. The latter is obtained below through an aggregated finite expansion and off-diagonal bounds uniform in the subsets of primes being tested.

The relevant classical inputs include zero-counting estimates, moments for the argument of zeta and for its increments, and mean-value estimates for Dirichlet polynomials. We retain the parameter restrictions of these results, including the admissible length of every product polynomial and the contribution of the boundary indices. The diagonal of an independent-prime model is transferred to the deterministic clock only after the nonzero frequencies are estimated.


\section{Notation and conventions}\label{sec:notation}
Put
\begin{equation}\label{eq:N0}
N_0(t)=\frac{t}{2\pi}\left(\log\frac{t}{2\pi}-1\right)+\frac78,
\qquad \lamb(t)=N_0'(t)=\frac1{2\pi}\log\frac{t}{2\pi},\qquad
N_0''(t)=\frac1{2\pi t}.
\end{equation}
For sufficiently large indices, $\tau_k$ is the unique solution $N_0(\tau_k)=k-\tfrac12$. Order the positive ordinates $\gamma_k$ of all nontrivial zeros with multiplicity. For $T>0$ let
\begin{equation}\label{eq:windows}
\calI=\{k:T<\gamma_k\le 2T\},\quad M_T=\#\calI,
\quad\av{F}=\frac1{M_T}\sum_{k\in\calI}F_k,\quad \norm{F}=\av{|F|^2}^{1/2}.
\end{equation}
We only use $T$ such that $M_T>0$ and all $\tau_k$ with $k\in\calI$ exceed $2\pi$.  The unconditional localization needed for these conventions is proved in \cref{lem:clocklocal}. The strict/weak endpoint convention in \eqref{eq:windows} is used in every sum.

For a real cutoff $X\ge2$, introduce the full Euler phase
\begin{equation}\label{eq:Sx}
S_X(t)=-\frac1\pi\sum_{p\le X}\sum_{m\ge1}
\frac{\sin(mt\log p)}{m p^{m/2}}.
\end{equation}
The inner series converges absolutely for each prime, and, since the outer sum is finite, termwise differentiation of any fixed order is legitimate. The first prime harmonic is $a_p\sin(t\log p)$ with $a_p=-1/(\pi\sqrt p)$. Define
\begin{equation}\label{eq:defs}
\begin{aligned}
s_k&=N_0(\tau_k)-N_0(\gamma_k),
&u_k(X)&=s_k-S_X(\gamma_k),\\
\mathfrak t_k(X)&=S_X(\gamma_k)-S_X(\tau_k),
&\tau_k^X&=\tau_k-\frac{S_X(\tau_k)}{\lamb(\tau_k)},\\
r_k(X)&=\lamb(\tau_k)(\tau_k^X-\gamma_k),
&d_k&=\tau_k-\gamma_k.
\end{aligned}
\end{equation}
All sampling points $\tau_k^X$ can be computed from indices and primes alone; the $\gamma_k$ appear only in the analysis of their errors.

\begin{longtable}{p{0.19\linewidth}p{0.75\linewidth}}
\toprule Symbol & Meaning and dependence\\\midrule\endfirsthead
\toprule Symbol & Meaning and dependence\\\midrule\endhead
$\rho_k,\beta_k,\gamma_k$ & Actual zero, real part, and positive ordinate; zeros counted with multiplicity.\\
$N_0,\lamb$ & Smooth count and its derivative; $\lamb(t)>0$ for $t>2\pi$.\\
$\tau_k$ & Midpoint clock, $N_0(\tau_k)=k-1/2$.\\
$T,\calI,M_T$ & Dyadic height, index window, exact cardinality.\\
$\langle\cdot\rangle_T,\|\cdot\|_{2,T}$ & Normalized average and its quadratic norm on $\calI$.\\
$X,Y$ & Finite prime cutoffs, potentially depending on $T$.\\
$S_X,a_p$ & Prime-power Euler phase and first-prime coefficient.\\
$s,u,\mathfrak t,r,d$ & Ordinal defect, prime residual, phase transport, normalized corrected error, bare-clock displacement.\\
$W_{X,Y}$ & Difference $S_X-S_Y$ between two finite cutoff phases.\\
$H_0,H_1,H_2$ & $\log\log T$, $\log H_0$, and $\log H_1$, defined for sufficiently large $T$.\\
$\omega,\theta,L$ & Auxiliary scale, cutoff exponent, and prime-band energy; explicitly defined where used.\\
$\mathcal J,\mu(\mathcal J)$ & A set of primes and its harmonic mass $\sum_{p\in\mathcal J}p^{-1}$.\\
$\mathcal R_p^\tau$ & Response coefficient of the clock sampled against the ordinal defect.\\
$\mathcal B_p$ & Response coefficient of an explicitly constructed phase difference.\\
$\Phi_T(p)$ & Characteristic factor of an independent-prime model, not a deterministic correlation by definition.\\
$u$ & Real Fourier frequency; $u=\log x$ in spectral prime-counting formulae.\\
$x_k^{(1)},e_{1,k}$ & First arithmetic correction at cutoff $X_1$, and its difference from the true ordinate.\\
$Z,X_1,Y,X_2$ & Two low-prime cutoffs, intermediate moving cutoff, and the second-clock cutoff.\\
$P_0,P_1,P_2,E_1$ & Finite prime-power decomposition and residual tail of the first correction.\\
$Q,Q^\flat,J_0,J_1$ & Second prime phase, its truncated polynomial, and the two Taylor orders.\\
$V_0,\theta_p$ & Variance scale of the prime phase, and moving-prime exponent $\log p/\log T$.\\
$\eps$ & Arbitrary strictly positive constant; its implicit-constant dependence must be tracked.\\
\bottomrule
\end{longtable}


\section{Published analytic inputs and their domains}\label{sec:inputs}
The following established external results enter the proofs below,
with the domains in which they are applied.
The Riemann--von Mangoldt formula for the full zero count is used
with the right-continuous convention at ordinates and with its smooth
remainder bounded by $O(t^{-1})$; the derivation of the cubic identity
itself is independent of the zeta function.

For the continuous prime approximation we use the fixed-order
Selberg--Fujii remainder estimate at $H=T$, with $a_2=1/5$ and
$k=2$ for the fourth moment, and at any fixed higher integer
$k$ used in a moment lemma.  The parameters remain in the original
range $T^{1/2+a_2}\le H\le T$, and no unbounded displacement
is inserted into Fujii's increment theorem.  The corrections
published in \cite{Fujii1981} are distinguished from the
continuous prime approximation employed here.

For a Dirichlet polynomial of length $R$, the Montgomery--Vaughan
mean-value theorem is used in the form \eqref{eq:MVinput}; for its
powers the applicable length is the product length, explicitly
tracked in \cref{lem:bandmom,lem:higherordinal,lem:polysampling,lem:aggerrors,lem:aggoffdiag}.  Kusmin--Landau is applied only when
the discrete phase derivative stays away from an integer; near the
product resonance we use the second-derivative test instead.
Hoeffding's inequality applies to the independent-prime model,
not by itself to a deterministic mean on the clock.  That transfer
is proved using the clock exponential-sum estimate.
The Hardy--Littlewood maximal-function theorem and its Sobolev
version \cite{Stein1970,Kinnunen1997} are used with the explicitly compactly supported derivative
of the prime phase, so that their $L^4$ and $L^8$ hypotheses are met.

The Gram-point moments of Korolev concern a distinct sampling system
and are not silently substituted for moments over the true-ordinate
window.  Similarly, the finite Euler-product construction in
\cite{Gonek2012} does not identify synthetic zeros with the
zeros of $\zeta$; it is pertinent as an independent analytic
Euler-product model, not as a spectral identification theorem.


\section{Exact accounting at finite height}\label{sec:exact}
\begin{lemma}[Uniform localization of the ordinal clock]
\label{lem:clocklocal}
There are constants $T_0,\Cloc>0$ such that, for every index with
$\gamma_k\ge T_0$, including all indices at multiple ordinates,
\begin{equation}\label{eq:clocklocal}
 |\tau_k-\gamma_k|\le\Cloc.
\end{equation}
Consequently, uniformly for $k\in\mathcal I_T$ and large $T$,
$\tau_k\in[T-\Cloc,2T+\Cloc]$ and
$\Lambda(\tau_k)\asymp\log T$.
\end{lemma}
\begin{proof}
Let $\gamma$ have multiplicity $m$ and let
$k=N(\gamma^-)+j$, $1\le j\le m$.  In the right- and
left-limit forms of the Riemann--von Mangoldt formula,
$N(t)=N_0(t)+S(t)+O(1/t)$, the standard bound
$S(t)=O(\log t)$ applies on both sides.  Their difference gives
$m=O(\log\gamma)$; hence
\[
 k-\frac12-N_0(\gamma)
 =N(\gamma^-)+j-\frac12-N_0(\gamma)
 =O(\log\gamma),
\]
with constants uniform in $j$ and in $m$.  Since
$N_0(\tau_k)=k-1/2$, this means
$|N_0(\tau_k)-N_0(\gamma_k)|\ll\log\gamma_k$.
For $\gamma$ sufficiently large, the increments
$N_0(\gamma+1)-N_0(\gamma)$ and
$N_0(\gamma)-N_0(\gamma-1)$ are both
$\asymp\log\gamma$ with positive coefficients.  Monotonicity
of $N_0$ for $t>2\pi$ then confines $\tau_k$ to
$[\gamma_k-C,\gamma_k+C]$ for a fixed sufficiently large absolute
$\Cloc$. Fix this same constant throughout the subsequent
window estimates; no step will assume $\Cloc\le1$.
This proves both assertions.
\end{proof}

\begin{lemma}[Taylor remainder with the true index window]\label{lem:taylor}
Suppose that $\gamma_k$ and $\tau_k$ belong to $[T/2,3T]$ for $k\in\calI$, with $T>4\pi$. Then
\begin{equation}\label{eq:taylore}
s_k=\lamb(\tau_k)d_k+e_k,\qquad
|e_k|\le\frac{d_k^2}{2\pi T}.
\end{equation}
Consequently, for every finite cutoff $X\ge2$,
\begin{equation}\label{eq:exact_r}
r_k(X)=u_k(X)+\mathfrak t_k(X)-e_k.
\end{equation}
\end{lemma}
\begin{proof}
Taylor's integral formula gives
$N_0(\gamma_k)=N_0(\tau_k)+N_0'(\tau_k)(\gamma_k-\tau_k)
+\int_{\tau_k}^{\gamma_k}(\gamma_k-v)N_0''(v)\,dv$.
Since $|N_0''(v)|\le1/(\pi T)$ on $[T/2,3T]$, the integral has absolute value at most $d_k^2/(2\pi T)$. Subtracting this identity from $N_0(\tau_k)$ proves the first assertion. By definition,
$r_k=\lamb(\tau_k)d_k-S_X(\tau_k)$, while
$u_k+\mathfrak t_k=s_k-S_X(\tau_k)$. This proves \eqref{eq:exact_r}.
\end{proof}

\begin{theorem}[Exact finite-height comparison]\label{thm:finite}
Under the assumptions of \cref{lem:taylor}, define $D_T=\max_{k\in\calI}|d_k|$. Then, for all finite cutoffs $X$, including a cutoff depending on $T$,
\begin{equation}\label{eq:norm_cmp}
\left|\norm{r(X)}-\norm{\mathfrak t(X)}\right|
\le\norm{u(X)}+\frac{D_T^2}{2\pi T},
\end{equation}
and
\begin{equation}\label{eq:square_cmp}
\left|\av{r(X)^2}-\av{\mathfrak t(X)^2}\right|
\le A_T(X)\bigl(2\norm{\mathfrak t(X)}+A_T(X)\bigr),
\quad A_T(X)=\norm{u(X)}+\frac{D_T^2}{2\pi T}.
\end{equation}
\end{theorem}
\begin{proof}
Take the $\ell^2(\calI)$ norm of \eqref{eq:exact_r}. The reverse triangle inequality bounds the difference of the two norms by $\norm{u-e}$, which is no greater than $A_T(X)$. For nonnegative numbers $a,b$, $|a^2-b^2|\le|a-b|(2b+|a-b|)$; apply this with the two norms.
\end{proof}

\begin{remark}[The position of the endpoints]
The window is defined by actual ordinates, not by $\tau_k$ or $\tau_k^X$. No replacement of $\{T<\gamma_k\le2T\}$ by $\{T<\tau_k\le2T\}$ is used in \cref{thm:finite}. Such a replacement, when necessary for a mean-value theorem, requires a separate bound for the indices in the symmetric difference. Merely knowing $|\tau_k-\gamma_k|\ll1$ does not make the boundary sets identical.
\end{remark}


\section{An exact cubic identity for the ordinal discrepancy}\label{sec:cubic}
The geometry of the counting function provides an exact identity which does
not require the phase-corrected count to be increasing.  The treatment of
multiple zeros is essential: replacing the full index set by distinct
ordinates changes the jump polynomials.

\begin{theorem}[Cubic counting identity]\label{thm:cubic}
Let $N(t)$ be a right-continuous integer-valued step function, with
locally finite jumps of positive integral multiplicity and indices
$k=1,2,\ldots$ attached in their natural order.  Let $F\in C^1[a,b]$,
where $a<b$, and let $\Gamma$ be the list of ordinates in $(a,b]$,
with multiplicity.  Put $G=N-F$ and $v_k=k-\tfrac12-F(\gamma_k)$.
Then, exactly,
\begin{equation}\label{eq:cubic}
\sum_{a<\gamma_k\le b}v_k^2
 =\int_a^b G(t)^2F'(t)\,dt
 +\frac{G(b)^3-G(a)^3}{3}-\frac{N(b)-N(a)}{12}.
\end{equation}
The identity holds with arbitrary jumps at either endpoint under the
right-continuous convention, and it does not require $F'\ge0$.
\end{theorem}
\begin{proof}
On an open interval containing no jump, $G'=-F'$ and hence
$(G^3)'=-3G^2F'$.  Consider a jump of size $m$ at $t_0$, and write
$g=G(t_0-)$.  The indices at this jump are
$N(t_0-)+1,\dots,N(t_0-)+m$; their discrepancies are
$g+j-\tfrac12$ ($1\le j\le m$).  Partition $[g,g+m]$ into its
unit subintervals.  The elementary identity
\[
\int_{j-1/2}^{j+1/2}z^2\,dz=j^2+\frac1{12}
\]
(valid after translating the integration variable) yields
\[
\frac{(g+m)^3-g^3}{3}
 =\sum_{j=1}^m(g+j-\tfrac12)^2+\frac{m}{12}.
\]
Summing the ordinary fundamental theorem of calculus between jumps
and this identity across every jump in $(a,b]$ gives
\[
G(b)^3-G(a)^3=-3\int_a^bG(t)^2F'(t)\,dt
 +3\sum_{a<\gamma_k\le b}v_k^2+\frac{N(b)-N(a)}4.
\]
This rearranges to \eqref{eq:cubic}.  All operations involve finitely
many jumps and integrals of bounded regulated functions on a compact
interval, so no interchange of limiting operations is involved.
\end{proof}

\begin{corollary}[The prime-corrected ordinal discrepancy]\label{cor:cubicprime}
Let $N(t)$ count all nontrivial zeta zeros with $0<\Im\rho\le t$,
with multiplicity, and set $F=N_0+S_X$.  With
$G_X=N-N_0-S_X$ and the conventions of \eqref{eq:windows},
\begin{equation}\label{eq:exactdefect}
M_T\av{u_k(X)^2}
 =\int_T^{2T}G_X(t)^2[\Lambda(t)+S'_X(t)]\,dt
 +\frac{G_X(2T)^3-G_X(T)^3}{3}-\frac{M_T}{12}.
\end{equation}
\end{corollary}
\begin{proof}
For each index $k$ one has
$v_k=k-\tfrac12-[N_0(\gamma_k)+S_X(\gamma_k)]=u_k(X)$.
Apply \cref{thm:cubic} on $(T,2T]$.  In particular a zero at $T$
is excluded and a zero at $2T$ is included, with its full multiplicity.
\end{proof}


\section{Continuous moments of prime phases}\label{sec:continuous}
Here and below an invocation of a published analytic theorem is distinguished
from its algebraic consequences.  The following form of the Montgomery--Vaughan
mean-value estimate is sufficient for our polynomial calculations:
for $A(t)=\sum_{n\le R}a_n n^{-it}$ and an interval $I$ of length $H$,
\begin{equation}\label{eq:MVinput}
\int_I|A(t)|^2dt=H\sum_{n\le R}|a_n|^2
 +O\left(\sum_{n\le R}n|a_n|^2\right).
\end{equation}
Its implied constant is absolute and its application to a product polynomial
requires using the \emph{length of that product}, not the length of either
factor.  For the prime reciprocal sum we use
$\sum_{p\le y}p^{-1}=\log\log y+B_1+O(1/\log y)$ for $y\ge3$.

\begin{lemma}[Prime-band moments]\label{lem:bandmom}
Let $2\le A\le B\le T^{1/2}$. Define
\[
P_{A,B}(t)=\sum_{A<p\le B}a_p\sin(t\log p),\qquad
V(A,B)=\sum_{A<p\le B}p^{-1}.
\]
Uniformly in this range,
\begin{align}
\int_T^{2T}P_{A,B}(t)^2dt
 &=\frac{T}{2\pi^2}V(A,B)+O(T^{1/2}+1),\label{eq:band2}\\
\int_T^{2T}|P_{A,B}(t)|^4dt
 &\ll T(1+V(A,B))^2.\label{eq:band4}
\end{align}
For the full prime-power phase and the same range of $B$,
\begin{equation}\label{eq:derivativeL2}
\int_T^{2T}|S'_B(t)|^2dt\ll T(1+\log B)^2.
\end{equation}
\end{lemma}
\begin{proof}
Put $D(t)=\sum_{A<p\le B}p^{-1/2-it}$.  The square of its imaginary
part equals $(|D|^2-\Re D^2)/2$.  Apply \eqref{eq:MVinput} to
$|D|^2$; the error is at most $\sum_{p\le B}p p^{-1}\le B$.
For the nonconjugated term, integrate each $(pq)^{-it}$ explicitly:
$|\int_T^{2T}(pq)^{-it}dt|\le2/\log4$.  Its total is at most
$C(\sum_{p\le B}p^{-1/2})^2\ll B$, using
$\sum_{n\le B}n^{-1/2}\ll\sqrt B$.  Hence the error in
\eqref{eq:band2} is $O(B)\subseteq O(T^{1/2})$.

For the fourth power use $|P_{A,B}|\le|D|/\pi$ and write
$D^2=\sum_{n\le B^2}b_nn^{-it}$.  By unique factorization,
$b_{p^2}=p^{-1}$ and $b_{pq}=2(pq)^{-1/2}$ when $p<q$.
Thus $\sum_n|b_n|^2\le2V(A,B)^2+\sum_pp^{-2}$.
As $B^2\le T$, the error $\sum nn|b_n|^2$ in
\eqref{eq:MVinput} is at most $T\sum|b_n|^2$.
This gives \eqref{eq:band4}.  For \eqref{eq:derivativeL2},
write the derivative as the sum of the first harmonic,
the squares, and all $m\ge3$ harmonics.  The latter series
is uniformly bounded because
$\sum_p\sum_{m\ge3}(\log p)p^{-m/2}<\infty$.
The first harmonic has coefficients $(\log p)/(\pi\sqrt p)$,
so \eqref{eq:MVinput}, partial summation and
$\sum_{p\le y}(\log p)/p=\log y+O(1)$ give
$O(T(1+\log B)^2)$.  The squares have coefficients
$(\log p)/(\pi p)$ and make a smaller contribution, including
the finite-length mean-value error.  Summing the three estimates
proves the lemma.
\end{proof}

\begin{lemma}[Moments of the unselected prime powers]\label{lem:primepowers}
For $B\le T^{1/2}$, write
$S_B=S_B^{\rm pr}+S_B^{\rm pp}$, where $S_B^{\rm pr}$ is the
first-harmonic sum.  On $[T,2T]$ one has
\begin{equation}\label{eq:pp2}
\int_T^{2T}|S_B^{\rm pp}(t)|^2dt\ll T,
\qquad \int_T^{2T}|S_B^{\rm pp}(t)|^4dt\ll T.
\end{equation}
\end{lemma}
\begin{proof}
The $m\ge3$ subseries is uniformly bounded.  The $m=2$ portion is
\[
-\frac1{2\pi}\Im\sum_{p\le B}p^{-1-2it}.
\]
Its second moment follows from \eqref{eq:MVinput} and the square
coefficients.  To estimate its fourth moment, square the polynomial.
The resulting frequencies are $(pq)^2$, with coefficient-square sum
bounded by a constant multiple of $(\sum_pp^{-2})^2$.
The mean-value error is bounded by
\[
C\sum_{p,q\le B}(pq)^2(pq)^{-2}
\ll B^2\le T.
\]
The uniform tail adds $O(T)$ to both bounds.
\end{proof}


\section{The ordinal prime-information threshold}\label{sec:threshold}
The analytical input needed in this section is a bounded fourth moment
of the remainder in the Selberg approximation, in the precise fixed-order,
continuous-height form stated next.  The reference to Fujii is to a
\emph{published} theorem, not to another manuscript of the author.

\begin{theorem}[Fixed-order Selberg--Fujii prime approximation]
\label{thm:selberginput}
Fix any integer $h\ge1$ and let
\begin{equation}\label{eq:fujiiorder}
 x_h=T^{1/(300h)},\qquad R_h=x_h^3=T^{1/(100h)},
 \qquad Y_h=T^{1/(80h)}.
\end{equation}
For sufficiently large $T$ depending on $h$,
\begin{equation}\label{eq:selbergfixedorder}
\int_T^{2T}\left|S(t)+\frac1\pi\sum_{p\le Y_h}
\frac{\sin(t\log p)}{\sqrt p}\right|^{2h}dt\ll_h T.
\end{equation}
In particular, for $h=2$, $Y_2=T^{1/160}$ and
\begin{equation}\label{eq:selberginput}
\int_T^{2T}\left|S(t)+\frac1\pi\sum_{p\le T^{1/160}}
\frac{\sin(t\log p)}{\sqrt p}\right|^4dt\ll T.
\end{equation}
This is a fixed-order consequence of the Selberg approximation
with Fujii's bounds on pp.~232--233 of \cite{Fujii1974},
with $a_2=1/5$, integration length $H=T$, and
auxiliary parameter $x_h=T^{a_2/(60h)}$.
The dependence of the constants on the \emph{fixed} integer
$h$ is retained: no growing-order uniformity is used in
this article. The domain $T^{1/2+a_2}\le H\le T$
is satisfied by $H=T$, and $x_h\ge2$ once
$T\ge2^{300h}$. The 1981 corrections to the increment
theorem of Fujii are not extrapolated to unbounded shifts.
\end{theorem}
\begin{proof}[Derivation and admissible polynomial lengths]
For fixed $h$, the Selberg formula printed on p.~232
of \cite{Fujii1974} represents $S(t)$ by the imaginary
part of a smoothed prime-power polynomial with
$n\le x_h^3=R_h$, together with five remainder terms.
The bounds on p.~233 yield, in the stated $H=T$ range,
\begin{equation}\label{eq:fujii5}
\int_T^{2T}|S(t)-S_{x_h,\mathrm{sm}}(t)|^{2h}dt
 \ll_h T.
\end{equation}
Indeed the $L^{2h}$ norms of the five errors are each
$O_h(T^{1/(2h)})$; their fixed-order sum is controlled
by Minkowski. The restrictions on Fujii's auxiliary
parameters are unchanged.

To pass to an ordinary first-prime sum, observe that the
smoothed coefficient agrees with the unsmoothed coefficient
for $p\le x_h$ and, for $x_h<p\le R_h$, its modulus is
bounded by $1/(\pi\sqrt p)$. Consequently the difference
of the first-prime parts of $S_{x_h,\mathrm{sm}}$ and
$S^{\rm pr}_{Y_h}$ has coefficients $O_h(p^{-1/2})$
supported on $x_h<p\le Y_h$. Its coefficient-square
sum is
\begin{equation}\label{eq:fujiiV}
\sum_{x_h<p\le Y_h}\frac1p
=\log\frac{\log Y_h}{\log x_h}+O(1)
=\log\frac{300}{80}+O(1)=O(1).
\end{equation}
If $D(t)$ is the corresponding complex prime polynomial,
$D(t)^h$ has length at most
$Y_h^h=T^{1/80}<T$. The Montgomery--Vaughan
mean-value theorem \eqref{eq:MVinput}, together with
unique factorization and the elementary multinomial
inequality, gives
\begin{equation}\label{eq:fujiiD}
\int_T^{2T}|D(t)|^{2h}dt
\ll_h (T+Y_h^h)
\left(\sum_{x_h<p\le Y_h}\frac1p\right)^h
\ll_h T.
\end{equation}
The prime-square terms in the smoothed polynomial have
frequencies $p^2\le R_h$, coefficients $O_h(p^{-1})$;
their $h$-th product polynomial has length at most
$R_h^h=T^{1/100}<T$, with bounded coefficient-square
sum for fixed $h$. All harmonics $p^m$ with $m\ge3$
have uniformly absolutely convergent coefficient series
because $\sum_p\sum_{m\ge3}p^{-m/2}/m<\infty$.
Thus their $2h$-th moments are $O_h(T)$ as well.
Minkowski applied to \eqref{eq:fujii5}, \eqref{eq:fujiiD}
and these prime-power bounds proves \eqref{eq:selbergfixedorder}.
Taking $h=2$ gives \eqref{eq:selberginput} exactly.
\end{proof}

\begin{lemma}[Extension of the fixed-order prime remainder]\label{lem:approxextension}
Under \cref{thm:selberginput}, uniformly for $2\le X\le T^{1/2}$,
with $c=\log X/\log T$ and $V=\sum_{\min(X,Y_0)<p\le\max(X,Y_0)}1/p$,
\begin{align}
\int_T^{2T}|S(t)-S_X(t)|^2dt
 &=\frac{T}{2\pi^2}V+O(T(1+\sqrt V)),\label{eq:EX2}\\
\int_T^{2T}|S(t)-S_X(t)|^4dt
 &\ll T(1+V)^2,\label{eq:EX4}\\
V&=\log(1/c)+O(1).\label{eq:Vsize}
\end{align}
The constants are absolute.
\end{lemma}
\begin{proof}
Write $S-S_X=E_0+W+P$, where $E_0$ is the remainder
in \eqref{eq:selberginput}, $W$ is the signed first-harmonic
prime band between $X$ and $Y_0$, and $P$ is a signed combination
of prime powers beyond first harmonics.  The latter has
uniformly bounded second and fourth moments by
\cref{lem:primepowers}; the $E_0$ second and fourth moments
are $O(T)$ by \eqref{eq:selberginput} and Holder.
\cref{lem:bandmom} gives $\int W^2=T V/(2\pi^2)+O(T^{1/2})$
and $\int W^4\ll T(1+V)^2$.  The cross term
$\int W(E_0+P)$ is $O(T\sqrt V+T)$ by Cauchy--Schwarz,
proving \eqref{eq:EX2}.  The elementary estimate
$|z_1+z_2+z_3|^4\le27\sum|z_i|^4$ proves
\eqref{eq:EX4}.

To see \eqref{eq:Vsize}, use the reciprocal-prime asymptotic
at the two endpoints.  If $X\le Y_0$, then
$V=\log\log Y_0-\log\log X+O(1)
 =\log(1/c)+\log(1/160)+O(1)$.
If $X>Y_0$, then $V\ll\log(80)+O(1)$, and
$\log(1/c)$ is bounded in that regime.
All statements are uniform, including the empty band and $X=Y_0$.
\end{proof}

\begin{theorem}[Unconditional mean-square ordinal prime defect]
\label{thm:primedefect}
There are absolute constants $C,T_0$ such that, for $T\ge T_0$
and $2\le X\le T^{1/2}$, with $c=(\log X)/(\log T)$,
\begin{equation}\label{eq:primedefect}
\left|\av{u_k(X)^2}-\frac1{2\pi^2}\log\frac1c\right|
 \le C\left(1+\sqrt{\log(1/c)}\right).
\end{equation}
The range includes $c=\frac12$ and the lower endpoint $X=2$;
the constants are uniform in $c$ in this entire interval.
\end{theorem}
\begin{proof}
We start from the exact identity \eqref{eq:exactdefect}.
The Riemann--von Mangoldt formula in the present normalization is
$N(t)=N_0(t)+S(t)+R(t)$, with $R(t)=O(1/t)$ at non-ordinate heights
(at an ordinate this is understood with the right-limit definition of $S$).
Write $E_X=S-S_X$.  The integrals in \eqref{eq:exactdefect}
are absolutely convergent over the finite interval, and
$G_X=E_X+R$.  By the classical bounds $S(t)=O(\log t)$ and
$\sum_{p\le X}p^{-1/2}\ll\sqrt X$, one has, uniformly
for $t\in[T,2T]$, $|G_X(t)|\ll\log T+\sqrt X\ll T^{1/4}$
for all sufficiently large $T$ (increasing the absolute constant
if necessary).  Consequently the cubic endpoint term is
$O(T^{3/4})$.  This argument includes heights where $N$ jumps;
we evaluate $N$ by its right limit, without shifting endpoints.

The contribution from replacing $G_X^2$ by $E_X^2$ in the integral
is bounded by
\[
O\left(T^{-1}\int_T^{2T}(|E_X|+1)(\Lambda+|S'_X|)dt\right).
\]
Cauchy--Schwarz, \eqref{eq:EX2} and \eqref{eq:derivativeL2}
give $O((1+V)^{1/2}\log T+\log T)$, a negligible term
once divided by $M_T\asymp T\log T$.

For the main integral, $\Lambda(t)=\Lambda(T)+O(1)$ for
$t\in[T,2T]$.  Hence, by \eqref{eq:EX2},
\begin{equation}\label{eq:thresholdmain}
\int_T^{2T}E_X(t)^2\Lambda(t)dt
=\frac{T\Lambda(T)}{2\pi^2}V
 +O(T\log T(1+\sqrt V)+T(1+V)).
\end{equation}
The derivative term is estimated by Cauchy--Schwarz:
\[
\left|\int_T^{2T}E_X(t)^2S'_X(t)dt\right|
\le\left(\int_T^{2T}E_X^4\right)^{1/2}
\left(\int_T^{2T}(S'_X)^2\right)^{1/2}
\ll T(1+V)(1+\log X).
\]
The Riemann--von Mangoldt count gives
$M_T=T\Lambda(T)+O(T)$.
Divide \eqref{eq:exactdefect} by $M_T$ and use these bounds.
The error is
\[
O\left(1+\sqrt V+\frac{1+V}{\log T}+(1+V)c
 +\frac{T^{-1/4}}{\log T}\right).
\]
Since $V\ll1+\log(1/c)$ and
$\sup_{0<c\le1/2}c\log(1/c)<\infty$,
this is $O(1+\sqrt{\log(1/c)})$ uniformly.
Finally replace $V$ by $\log(1/c)+O(1)$ using
\eqref{eq:Vsize}.  This completes the proof.
\end{proof}



\section{Clock exponential sums and endpoint orthogonality}\label{sec:clockorth}
For this section $T$ is sufficiently large, the indices are those of
\eqref{eq:windows}, and the inverse $x\mapsto\tau_x$ is defined by
$N_0(\tau_x)=x-\tfrac12$.  The half-integer convention fixes every phase.

\begin{lemma}[Kusmin--Landau sum on the exact ordinal interval]
\label{lem:klclock}
If $0<|w|\le\tfrac34\log T$, then for every consecutive subinterval
$J\subset\mathcal I_T$,
\begin{equation}\label{eq:klclock}
\left|\sum_{k\in J}e^{iw\tau_k}\right|
 \ll1+\frac{\log T}{|w|}.
\end{equation}
More generally, if the complex weights $z_k$ have supremum norm $A_0$
and total variation $A_1$, the weighted sum is bounded by
$O((A_0+A_1)(1+\log T/|w|))$.
\end{lemma}
\begin{proof}
The derivative of $f(x)=w\tau_x/(2\pi)$ is
$f'(x)=w/\log(\tau_x/(2\pi))$.  On the relevant interval,
$\tau_x\in[T-O(1),2T+O(1)]$, so $|f'|\le0.8$ for large $T$;
$f'$ is monotone, and its distance from the closest integer is at least
$\min(|w|/\log(3T),1/5)$.
The Kusmin--Landau first-derivative estimate for exponential sums
therefore bounds all consecutive partial sums by
$O(1+\log T/|w|)$.  Abel summation proves the weighted assertion,
including weights that vary with the index.  The definition
$(T,2T]$ makes the sampled indices consecutive even if an endpoint
is a multiple ordinate.
\end{proof}

\begin{lemma}[Almost orthogonal first prime harmonics]\label{lem:orthprimes}
For $0<b\le\tfrac12$ let
$v_p(k)=\sqrt{2/M_T}\sin(\tau_k\log p)$ for primes $p\le T^b$,
and let $G=(\sum_kv_p(k)v_q(k))_{p,q}$.  Then
\begin{equation}\label{eq:orthprimes}
\|G-I\|_{\mathrm{op}}
 \ll T^{b-1}\log T+\frac1{\log T}.
\end{equation}
The estimate remains valid for every subset of those primes, with the
same absolute implicit constant, including the endpoint $b=1/2$.
\end{lemma}
\begin{proof}
Set $C_T(y)=M_T^{-1}\sum_{k\in\mathcal I_T}e^{i\tau_k\log y}$.
By the product identity for sines,
$G_{pq}=\Re C_T(p/q)-\Re C_T(pq)$.  For $p\ne q$, the first
factor satisfies
\[
|C_T(p/q)|\ll\frac{\max(p,q)}{T|p-q|}
\]
by \cref{lem:klclock} and $M_T\asymp T\log T$.
The sum of the absolute entries of any row off the diagonal is
$O(T^{b-1}\log T)$; Schur's test gives the corresponding
operator-norm bound.

The product terms $C_T(pq)$ require separate treatment near the
integer derivative $f'=1$; applying the first-derivative test there
would be invalid.  Put $n=pq\le T$.  If $n\le T^{3/4}$,
\cref{lem:klclock} gives $C_T(n)\ll T^{-1}$.  If
$T^{3/4}<n\le T/\log T$, then
\[
\frac12\le\frac{\log n}{\log(\tau_x/2\pi)}
 \le 1-\frac{\log\log T}{2\log T}
\]
for $T$ large, so the first-derivative estimate still gives
$C_T(n)\ll T^{-1}$.  In the remaining range
$T/\log T<n\le T$, the second derivative of
$f(x)=\tau_x\log n/(2\pi)$ has absolute size
$\asymp1/(T\log^2T)$ and bounded ratio of its maximum to minimum.
Van der Corput's second-derivative estimate yields
$|\sum_{k\in\mathcal I_T}e^{i\tau_k\log n}|
\ll\sqrt T\log T$, whence $C_T(n)\ll T^{-1/2}$.
Every product in this exceptional last range has
$p,q>\sqrt T/\log T$ because $p,q\le\sqrt T$.
The number of such pairs is $O(T/\log^2T)$, using the usual
upper bound for $\pi(\sqrt T)$.  Consequently the Hilbert--Schmidt
norm of the product matrix satisfies
\[
\sum_{p,q\le T^b}|C_T(pq)|^2
 \ll \frac{T}{T^2}+\frac{T}{\log^2T}\frac1T
 \ll(\log T)^{-2}.
\]
The operator norm is at most its Hilbert--Schmidt norm.  Combining
with the quotient matrix proves \eqref{eq:orthprimes}.
\end{proof}

\begin{corollary}[Harmonic mass on arbitrary bands]
\label{cor:massbands}
For every prime subset $\mathcal J\subseteq\{p:p\le T^{1/2}\}$,
\begin{equation}\label{eq:massbands}
\av{W_{\mathcal J}(\tau_k)^2}
=\frac1{2\pi^2}\sum_{p\in\mathcal J}\frac1p
 \left(1+O((\log T)^{-1})\right).
\end{equation}
The error is interpreted as zero if $\mathcal J$ is empty.
\end{corollary}
\begin{proof}
Use the Gram matrix of \cref{lem:orthprimes} and the coefficient
vector $a_p=-1/(\pi\sqrt p)$ supported on $\mathcal J$.
The scalar quadratic form differs from its identity-matrix value by
at most the operator norm times the coefficient-square sum.
\end{proof}


\section{Sampling of corrected points}\label{sec:sampling}
The first correction changes the sampling set and can produce
multiple points in a short interval, although the original clock is
separated.  The relevant multiplicity is quadratic, not merely the
maximum number of points in a cell.

\begin{lemma}[Finite-point sampling inequality]\label{lem:abstractsampling}
Let $y_1,\ldots,y_M$ be real points, $b>0$, and
$n_b(t)=\#\{j:t-b\le y_j\le t\}$.  For every absolutely
continuous nonnegative function $F$ on an interval containing
all $[y_j,y_j+b]$,
\begin{equation}\label{eq:abssampling}
\sum_{j=1}^MF(y_j)
\le\frac{\|n_b\|_2}{b}
\left(\|F\|_2+b\|F'\|_2\right),
\end{equation}
where all $L^2$ norms on the right are over that interval.
\end{lemma}
\begin{proof}
For every $z\in[y_j,y_j+b]$, absolute continuity gives
$F(y_j)\le F(z)+\int_{y_j}^{y_j+b}|F'(v)|dv$.
Average over $z$ and sum over $j$.  The first contribution is
$b^{-1}\int n_b(z)F(z)dz$; the second is bounded by
$\int n_b(z)|F'(z)|dz$.  Apply Cauchy--Schwarz to both
integrals.  The definition of $n_b$ makes this estimate exact,
regardless of coincidence among the points.
\end{proof}

\begin{lemma}[Moment sampling with an additive error]\label{lem:polysampling}
Let $f$ be a real trigonometric polynomial supported on prime powers
$n=p^m\le R$, with no constant term, and coefficients $f_n$.
Define
$\|f\|_1=\sum_n|f_n|$ (counting the sine and cosine coefficients),
and
\[
\sigma_f^2=\sum_p\left(\sum_{p^m\le R}|f_{p^m}|\right)^2.
\]
Let $q\in\{1,2,3,4\}$, let $y_k\in[T-T^{1/2},2T+T^{1/2}]$,
and suppose that, for $b=2\pi/\log(3T)$,
\begin{equation}\label{eq:samplingconditions}
\int_{\mathbb R}n_b(t)^2dt\le\varkappa bM_T,
\qquad (2\|f\|_1)^{4q}R^{8q}\le T^{1/2},
\end{equation}
where $\varkappa\ge1$.  Then
\begin{equation}\label{eq:polysampling}
\av{|f(y_k)|^{2q}}
\ll_q\varkappa^{1/2}(1+b\log R)
\left(\sigma_f^{2q}+T^{-1/4}\right).
\end{equation}
The additive $T^{-1/4}$ term cannot be absorbed into
$\sigma_f^{2q}$ without a lower bound for that variance.
\end{lemma}
\begin{proof}
Let $F=f^{2q}$.  Apply \cref{lem:abstractsampling} on
$J=[T-2T^{1/2},2T+2T^{1/2}]$.
In the independent-prime model, the contributions of distinct primes
are independent, centered and bounded by their coefficient sums;
Hoeffding's inequality, integrated against $2qz^{2q-1}dz$,
gives $\mathbb E|f|^{4q}\ll_q\sigma_f^{4q}$.
The rational trigonometric polynomial $f^{4q}$ has height at most
$R^{4q}$ and coefficient norm at most $(2\|f\|_1)^{4q}$.
For a nonzero rational frequency $\log(a/b)$ with $a,b\le R^{4q}$,
$|\int_Je^{it\log(a/b)}dt|\le2R^{4q}$.
The model zero-frequency term is its expectation by unique
factorization.  Therefore
\[
\frac1{|J|}\int_J|f|^{4q}dt
\ll_q\sigma_f^{4q}+\Delta_f,
\qquad \Delta_f=\frac{(2\|f\|_1)^{4q}R^{4q}}T.
\]
The same argument for $f'$ gives the bound with an additional
$(\log R)^{4q}$, since differentiation multiplies each coefficient
by at most $\log R$.  Holder's inequality applied to
$F'=2q f^{2q-1}f'$ now yields
$\|F\|_2+b\|F'\|_2
\ll_q\sqrt{|J|}(1+b\log R)(\sigma_f^{4q}+\Delta_f)^{1/2}$.
Use $|J|\asymp bM_T\asymp T$ in \eqref{eq:abssampling}.
Finally the second hypothesis of \eqref{eq:samplingconditions}
shows $\Delta_f\le T^{-1/2}$, proving \eqref{eq:polysampling}.
\end{proof}


\begin{lemma}[Maximal derivatives at deformed sampling points]
\label{lem:maxsample}
Let $f$ be a real, finite prime-power polynomial with frequencies
at most $\log R$, whose eighth-power product length and coefficient
norm satisfy the mean-value condition of
\cref{lem:polysampling} with $q=2$. Let $\varphi$ be supported in
$[T-3T^{1/2},2T+3T^{1/2}]$, equal to one on
$[T-2T^{1/2},2T+2T^{1/2}]$, with derivatives bounded
by absolute constants. Set $g=\mathcal M(\varphi f)$, where
$\mathcal M$ is the centered Hardy--Littlewood maximal operator.
Suppose the points $y_k$ have the multiplicity bound
$\int n_b^2\le CbM_T$ and lie in the inner interval. Then
\begin{equation}\label{eq:maxsample}
\bigl\langle |g(y_k)|^4\bigr\rangle_T^{1/4}
\ll_C (1+b\log R)\left(\sigma_f+T^{-1/16}\right),
\end{equation}
where $\sigma_f$ is the prime-group coefficient variance proxy
of \cref{lem:polysampling}.  This applies with
$f=(Q^\flat)'$ and gives
\begin{equation}\label{eq:sampleQprime}
\bigl\langle|\mathcal M(\varphi(Q^\flat)')(y_k)|^4
\bigr\rangle_T^{1/4}\ll 1+\log X_2.
\end{equation}
\end{lemma}
\begin{proof}
By the strong $L^8$ bound for the maximal operator and the
Sobolev maximal theorem (in the form $|g'|\le
\mathcal M((\varphi f)')$ a.e.),
\begin{equation}\label{eq:g8}
\|g\|_{L^8}\ll\|\varphi f\|_{L^8},\qquad
\|g'\|_{L^8}\ll\|(\varphi f)'\|_{L^8}.
\end{equation}
The latter derivative has the terms $\varphi' f$ and
$\varphi f'$.  Expanding $f^8$ into its finite rational
frequencies and applying the clock-independent time averaging
argument in the proof of \cref{lem:polysampling}, with $q=2$,
gives, on every interval of length comparable to $T$,
\[
\int |f(t)|^8dt\ll T(\sigma_f^8+T^{-1/2}),\quad
\int|f'(t)|^8dt\ll
T(\log R)^8(\sigma_f^8+T^{-1/2}).
\]
The cutoff terms are treated identically by domination and its
bounded derivatives. Thus
\[
\|g\|_{L^8}+b\|g'\|_{L^8}
\ll T^{1/8}(1+b\log R)
(\sigma_f+T^{-1/16}).
\]
Apply \cref{lem:abstractsampling} to $F=g^4$.
Since $(g^4)'=4g^3g'$ a.e., H\"older gives
\[
\|F\|_2\le\|g\|_8^4,\qquad
\|F'\|_2\le4\|g\|_8^3\|g'\|_8.
\]
The multiplicity factor, divided by $M_T$, is
$\sqrt{CbM_T}/(bM_T)\asymp T^{-1/2}$, and
$\|g\|_8^4=O(T^{1/2})$ times the fourth power of the
coefficient scale. This proves \eqref{eq:maxsample}.
For $f=(Q^\flat)'$, coefficients on the prime $p$ have size
$\ll(\log p)/\sqrt p$ and its higher-power companions; hence
\[
\sigma_f^2\ll\sum_{p\le X_2}\frac{(\log p)^2}{p}
\ll(1+\log X_2)^2.
\]
The length requirement holds because $R=X_2^2=T^{o(1)}$
while the eighth product and the coefficient norm are
$T^{o(1)}$. Moreover $b\log R=O(H_1^{-1/2})$.
This completes \eqref{eq:sampleQprime}.
\end{proof}

\begin{lemma}[Quadratic multiplicity for first-step deformations]
\label{lem:multcorrected}
Let $X_1=T^{1/\sqrt{H_0}}$, $b=2\pi/\log(3T)$,
and let $\Pi(t)$ be a finite coefficientwise contraction of the
full $S_{X_1}(t)$ (including its absolutely convergent prime-power
tails).  Set
\[
y_k=\tau_k-\frac{\Pi(\tau_k)}{\Lambda(\tau_k)},\qquad
n_b(t)=\#\{k\in\mathcal I_T:t-b\le y_k\le t\}.
\]
Then $y_k\in[T-T^{1/2},2T+T^{1/2}]$ and
\begin{equation}\label{eq:multcorrected}
\int n_b(t)^2dt\ll bM_T
\end{equation}
uniformly in every such contraction, including
$\Pi=P_0+vP_1$, $\Pi=vP_0$, and
$\Pi=P_0+P_1+v(P_2+E_1)$, $v\in[0,1]$.
\end{lemma}
\begin{proof}
For $T$ sufficiently large, the crude bound
$\sup|\Pi|\ll X_1$ gives
$|y_k-\tau_k|\ll X_1/\log T<T^{1/2}/2$.
The derivative satisfies
\begin{equation}\label{eq:PiL2}
\int_{T-\Cloc-2}^{2T+\Cloc+2}|\Pi'(t)|^2dt
\ll T(1+\log X_1)^2,
\end{equation}
by the same Montgomery--Vaughan calculation as
\eqref{eq:derivativeL2}, uniformly under coefficient domination.
The enlarged interval has length $T+2\Cloc+4\asymp T$;
its shift and its bounded extra length do not change the
mean-value bound. By \cref{lem:clocklocal}, it contains all
clock points $\tau_k$ for $k\in\mathcal I_T$.
The fourth moment at the clock obeys
\begin{equation}\label{eq:PiL4}
\av{|\Pi(\tau_k)|^4}\ll(1+H_0)^2.
\end{equation}
Indeed, truncate the full phase to prime powers $\le X_1^2$;
the tail is $O(1/\log X_1)$ uniformly.  The fourth power of the
remaining polynomial has frequency height $\le X_1^8=T^{8/\sqrt{H_0}}
=T^{o(1)}$.  Apply \cref{lem:klclock} to all nonzero frequencies
and identify the diagonal with the independent-prime fourth moment.
The latter is bounded by the Hoeffding inequality, since
$\sum_{p\le X_1}p^{-1}\ll H_0$.

We compute the second moment of the multiplicity function exactly:
\begin{equation}\label{eq:paircount}
\int n_b(t)^2dt
=\sum_{j,k\in\mathcal I_T}(b-|y_j-y_k|)_+.
\end{equation}
The diagonal and $|j-k|\le1$ pairs contribute at most $3bM_T$.
Fix $m\ge2$ and a pair $(k,k+m)$ counted in
\eqref{eq:paircount}.  Since the clock gaps are at least $b$,
$\tau_{k+m}-\tau_k\ge mb$, and hence the shift
$h(t)=\Pi(t)/\Lambda(t)$ satisfies
$|h(\tau_{k+m})-h(\tau_k)|\ge(m-1)b$.
By the derivative formula for $h$ and $\Lambda'/\Lambda^2
\ll (T\log^2T)^{-1}$, its contribution involving $\Pi$
itself is $o((m-1)b)$; thus
$\int_{\tau_k}^{\tau_{k+m}}|\Pi'(t)|dt\gg(m-1)b\log T$.
The interval length is at most $mb(1+O(1/\log T))$.
Cauchy--Schwarz shows that its $L^2$ derivative integral is
$\gg mb(\log T)^2$ for $m\ge2$.
The intervals associated to a fixed $m$ overlap at most $m$ times;
using \eqref{eq:PiL2} gives
\begin{equation}\label{eq:pairvar}
\#\{k:|y_{k+m}-y_k|<b\}
 \ll \frac{T(1+\log X_1)^2}{b(\log T)^2}
 \ll M_T H_0^{-1}.
\end{equation}
Alternatively, at least one of the shifts has magnitude
$(m-1)b/2$, so the corresponding $|\Pi(\tau)|\gg m-1$.
Markov's inequality applied to \eqref{eq:PiL4} gives
\begin{equation}\label{eq:pairsiz}
\#\{k:|y_{k+m}-y_k|<b\}
 \ll \frac{M_T(1+H_0)^2}{(m-1)^4}.
\end{equation}
Sum the minimum of \eqref{eq:pairvar} and \eqref{eq:pairsiz}
over $m\ge2$.  Splitting where $m\asymp H_0^{3/4}$,
this sum is $O(M_TH_0^{-1/4})$.
Combine with \eqref{eq:paircount} to conclude
$\int n_b^2\le (3+O(H_0^{-1/4}))bM_T$.
The argument depends only on coefficientwise bounds, so the constants
are uniform in the indicated deformation parameters.
\end{proof}



\section{The second sampling scale and the residual}\label{sec:second}
The first corrected point is denoted $x_k^{(1)}$, avoiding the use of
$\rho_k$ for anything other than a genuine zero.
For sufficiently large $T$ put
\begin{equation}\label{eq:scales}
\begin{gathered}
H_0=\log\log T,\quad H_1=\log H_0,\quad H_2=\log H_1,\\
X_1=T^{H_0^{-1/2}},\qquad X_2=T^{H_1^{-1/2}},\\
x_k^{(1)}=\tau_k-\frac{S_{X_1}(\tau_k)}{\Lambda(\tau_k)}.
\end{gathered}
\end{equation}
The independent arithmetic definition uses no actual ordinate.

\begin{lemma}[Higher fixed moments at actual ordinates]\label{lem:higherordinal}
For every fixed integer $m\ge1$, provided the prime cutoff
$X=T^c$ satisfies $c\to0$ and $c\le1/(16m)$, one has
\begin{equation}\label{eq:higherordinal}
\av{|s_k-S_X(\gamma_k)|^{2m}}
 \ll_m\left(1+\log\frac1c\right)^m.
\end{equation}
In particular, for $X=X_1$ and all sufficiently large $T$,
$\|u(X_1)\|_{4,T}\ll\sqrt{H_1}$, whereas
$\|s\|_{8,T}\ll H_0^{1/2}$.
\end{lemma}
\begin{proof}
Let $B_n(z)$ be the Bernoulli polynomials and set
$\mathcal P_m(z)=B_{2m+1}(z+\frac12)/(2m+1)$.
The Bernoulli difference identity
$B_n(z+1)-B_n(z)=nz^{n-1}$ gives
\[
\mathcal P_m(z+r)-\mathcal P_m(z)
 =\sum_{j=1}^r(z+j-\tfrac12)^{2m}
\]
for every positive integer $r$.
Repeat the proof of \cref{thm:cubic} with $\mathcal P_m(G)$.
The derivative $\mathcal P_m'(z)$ is an even polynomial of degree
$2m$, with leading term $z^{2m}$ and lower even powers.
This gives an exact sum over $(T,2T]$ expressed by integrals of
$\mathcal P_m'(N-N_0-S_X)(\Lambda+S_X')$ and its endpoints.

For the required continuous $4m$-th moment, apply
\cref{thm:selberginput} with its order parameter $h=2m$:
\[
Y_{2m}=T^{1/(160m)},\qquad
\int_T^{2T}|S-S_{Y_{2m}}^{\rm pr}|^{4m}dt\ll_m T.
\]
This is exactly the original fixed-order Selberg--Fujii domain,
with $a_2=1/5$, $H=T$ and
$x_{2m}=T^{1/(600m)}\ge2$ once $T$ is sufficiently large.
The first-prime difference between $Y_{2m}$ and $X=T^c$
is a signed prime polynomial. Its $2m$-th power has
product length at most
$\max(X,Y_{2m})^{2m}\le T^{1/8}$,
because $c\le1/(16m)$ and
$Y_{2m}^{2m}=T^{1/80}$. Its coefficient-square
variance is at most
\[
\sum_{\min(X,Y_{2m})<p\le\max(X,Y_{2m})}\frac1p
\ll_m 1+\log\frac1c.
\]
The Montgomery--Vaughan mean-value theorem applied to
that product polynomial yields the $4m$-th moment bound
$\ll_m T(1+\log(1/c))^{2m}$. For the omitted
prime powers, the square-prime part has coefficients $O(p^{-1})$
and product length $\le X^{4m}\le T^{1/4}$;
its coefficient-square sum is bounded for fixed $m$.
The harmonics $p^j$ with $j\ge3$ form a uniformly
absolutely convergent series. Thus the $4m$-th moment
of the full prime-power difference is bounded by the
same quantity. Minkowski and Holder give, with
$E_X=S-S_X$,
\[
\int_T^{2T}|E_X|^{4m}dt\ll_m
 T(1+\log(1/c))^{2m},\qquad
\int_T^{2T}|E_X|^{2m}dt\ll_m
 T(1+\log(1/c))^{m}.
\]

The derivative term has absolute value at most
\[
\left(\int_T^{2T}|E_X|^{4m}dt\right)^{1/2}
\left(\int_T^{2T}|S_X'|^2dt\right)^{1/2}
\ll_m T(1+\log(1/c))^m(1+\log X).
\]
After dividing by $M_T\asymp T\log T$ this is
$O_m((1+\log(1/c))^m)$, because $c\le1/(16m)$.
The lower-degree Bernoulli terms satisfy the same bound.
The endpoint polynomial has degree $2m+1$ and, by
$|N-N_0-S_X|\ll\log T+\sqrt X$, contributes
$O_m((\log T+\sqrt X)^{2m+1})$; since
$c(2m+1)/2<1$, it is $o(M_T)$ times any positive fixed
power of $\log\log T$.  The smooth Riemann--von Mangoldt
remainder is handled as in \cref{thm:primedefect}.
This proves \eqref{eq:higherordinal}.
For $s$, the continuous eighth moment is derived from
\cref{thm:selberginput} with $h=4$:
$Y_4=T^{1/320}$ and
$\int_T^{2T}|S-S_{Y_4}^{\rm pr}|^8\ll T$.
The fourth power of the prime polynomial has length
$Y_4^4=T^{1/80}<T$ and variance
$\sum_{p\le Y_4}p^{-1}\ll H_0$;
\eqref{eq:MVinput} then gives
$\int_T^{2T}|S_{Y_4}^{\rm pr}|^8dt\ll T H_0^4$.
Thus $\int_T^{2T}|S(t)|^8dt\ll T H_0^4$,
including the correct fixed-order input and its range.
Repeat the Bernoulli counting identity with $F=N_0$:
the main integral uses this eighth moment,
its lower even powers are smaller, and its endpoint
is $O((\log T)^9)=o(T\log T)$. Since $F'=\Lambda$,
this proves $\av{|s|^8}\ll H_0^4$ without invoking
any unstated external eighth-moment formula.
\end{proof}

\begin{lemma}[The short first transport]\label{lem:shorttransport}
For $2\le X\le T^{1/2}$,
\begin{equation}\label{eq:shorttransport}
\left\|S_X(\gamma_k)-S_X(\tau_k)\right\|_{2,T}
 \ll\sqrt{H_0}\,\frac{1+\log X}{\log T}.
\end{equation}
\end{lemma}
\begin{proof}
The clock displacement satisfies
$\|\gamma-\tau\|_{4,T}\ll\sqrt{H_0}/\log T$
by the fourth-moment case of \cref{lem:higherordinal} at $X=2$
and Taylor's formula for $N_0$.
Let $g(t)$ be the centered Hardy--Littlewood maximal function of
$\varphi(t)S'_X(t)$, where $\varphi$ is a smooth cutoff equal to
one on $[T-\Cloc-1,2T+\Cloc+1]$ and supported in
$[T-\Cloc-3,2T+\Cloc+3]$, with $\Cloc$ as in
\cref{lem:clocklocal}. Its first derivative is uniformly bounded.
Every segment joining $\tau_k$ to $\gamma_k$ lies in the
inner interval. Hence the maximal-function integral below is
the integral of the \emph{full} derivative along that segment.
The integral from $\tau_k$ to $\gamma_k$ obeys
$|S_X(\gamma_k)-S_X(\tau_k)|
\le2|\gamma_k-\tau_k|g(\tau_k)$.
The $L^4$ boundedness of this maximal operator, together with its
Sobolev derivative bound and the fourth moments of $S_X'$ and
$S_X''$ obtained from \eqref{eq:MVinput} (raise the prime
polynomials to their squares before taking the mean), gives
\[
\int_{\mathbb R}g^4\ll T(1+\log X)^4,\qquad
\int_{\mathbb R}|g'|^4\ll T(1+\log X)^8.
\]
Use the disjoint cells $[\tau_k,\tau_k+b]$, with
$b=2\pi/\log(3T)$, and one-dimensional Sobolev averaging to obtain
$\av{g(\tau_k)^4}\ll(1+\log X)^4$.
Cauchy--Schwarz and the fourth moment of the clock displacement
establish \eqref{eq:shorttransport}.
\end{proof}


\begin{lemma}[Fourth-moment transport at the first arithmetic scale]
\label{lem:firstL4}
For $X_1=T^{H_0^{-1/2}}$ and sufficiently large $T$,
\begin{equation}\label{eq:firstL4}
\bigl\|S_{X_1}(\gamma)-S_{X_1}(\tau)\bigr\|_{4,T}
\ll 1,
\qquad\bigl\|x^{(1)}-\gamma\bigr\|_{4,T}
\ll\frac{\sqrt{H_1}}{\log T}.
\end{equation}
\end{lemma}
\begin{proof}
By \cref{lem:higherordinal} with $s$ at exponent eight,
\[
\|\gamma-\tau\|_{8,T}\ll\frac{\sqrt{H_0}}{\log T},
\]
using the bounded displacement and the derivative
$\Lambda(v)\asymp\log T$ on the segment from each ordinate to
its clock point.  With the compactly supported cutoff from the
proof of \cref{lem:shorttransport}, let
$g=\mathcal M(\varphi S'_{X_1})$.
The enlarged supports are contained in an interval of length
$T+O(\Cloc)$, so the same mean-value theorem applies
with absolute constants. The $8$th-power mean of the prime
part of $S'_{X_1}$ follows
by applying Montgomery--Vaughan to the fourth power of its
Dirichlet polynomial: its product length is at most
$X_1^4=T^{4/\sqrt{H_0}}=T^{o(1)}$.
Its coefficient-square sum is bounded by
$C(\sum_{p\le X_1}\log^2p/p)^4\ll(\log X_1)^8$.
The second harmonics are treated with frequencies $p^2$ and
product length $X_1^8=T^{o(1)}$; higher harmonics and their
first two derivatives have absolutely convergent tails.
Consequently
\begin{equation}\label{eq:continuous8}
\int_{T-3}^{2T+3}|S'_{X_1}(t)|^8dt
 \ll T(1+\log X_1)^8,
\quad
\int_{T-3}^{2T+3}|S''_{X_1}(t)|^8dt
 \ll T(1+\log X_1)^{16}.
\end{equation}
The strong $L^8$ bound for the centered maximal operator,
together with its Sobolev derivative estimate, now gives
$\int g^8\ll T(1+\log X_1)^8$ and
$\int|g'|^8\ll T(1+\log X_1)^{16}$.
Since the disjoint clock cells each have length
$b=2\pi/\log(3T)$, averaging $g^8$ over those cells and
using absolute continuity yields
\[
\av{g(\tau_k)^8}^{1/8}
 \ll (1+b\log X_1)(1+\log X_1)
 \ll 1+\log X_1.
\]
The pointwise integral estimate
$|S_{X_1}(\gamma_k)-S_{X_1}(\tau_k)|
\le2|\gamma_k-\tau_k|g(\tau_k)$ and H\"older
therefore give
\[
\bigl\|S_{X_1}(\gamma)-S_{X_1}(\tau)\bigr\|_{4,T}
 \ll \frac{\sqrt{H_0}}{\log T}(1+\log X_1)=O(1).
\]
Finally \eqref{eq:exact_r} gives
$\Lambda(\tau)(x^{(1)}-\gamma)=
u(X_1)+
[S_{X_1}(\gamma)-S_{X_1}(\tau)]-e$.
By \cref{lem:higherordinal}, $\|u(X_1)\|_{4,T}
\ll\sqrt{H_1}$, and $\|e\|_{4,T}\ll T^{-1}$ by the
bounded displacement. Divide by $\Lambda(\tau)\asymp\log T$
to finish.
\end{proof}

\begin{lemma}[Second-scale replacement of the actual ordinates]
\label{lem:secondreplacement}
With $Q=S_{X_2}$ and the first corrected points in
\eqref{eq:scales},
\begin{equation}\label{eq:secondreplacement}
\left\|s_k-Q(x_k^{(1)})\right\|_{2,T}
\ll1+\sqrt{H_2}.
\end{equation}
\end{lemma}
\begin{proof}
The identity
$s_k-Q(x_k^{(1)})=u_k(X_2)+Q(\gamma_k)-Q(x_k^{(1)})$
is exact.  By \cref{thm:primedefect}, the first term has norm
$O(1+\sqrt{H_2})$.
Let $e_{1,k}=x_k^{(1)}-\gamma_k$.
By \cref{lem:firstL4},
$\|e_1\|_{4,T}\ll\sqrt{H_1}/\log T$,
with all eighth-moment polynomial lengths checked at $X_1=T^{o(1)}$.

To control the full Euler phase, write $Q=Q^\flat+H_2^\flat$, where
\begin{equation}\label{eq:Qflatlocal}
Q^\flat(t)=\sum_{p\le X_2}\sum_{p^m\le X_2^2}
\beta_{p^m}\sin(mt\log p),\qquad
H_2^\flat=Q-Q^\flat.
\end{equation}
For every $R\ge2$, the omitted $m$ for each $p\le R$ begin
strictly after $m>2\log R/\log p$, so the geometric series gives
\begin{equation}\label{eq:derivativetail}
\begin{split}
\|S_R-S_R^\flat\|_\infty
&\le \frac{C}{R}\sum_{p\le R}\frac1{1-p^{-1/2}}
 \ll \frac1{\log R},\\
\|(S_R-S_R^\flat)'\|_\infty
&\le\frac{C}{R}\sum_{p\le R}
 \frac{\log p}{1-p^{-1/2}}\ll1.
\end{split}
\end{equation}
Here the first estimate uses $\pi(R)\ll R/\log R$,
and the second $\sum_{p\le R}\log p\ll R$;
the constants are absolute. Absolute and uniform convergence
justifies termwise differentiation for each finite prime cutoff.
Therefore
\begin{equation}\label{eq:tailtransport}
\|H_2^\flat(\gamma)-H_2^\flat(x^{(1)})\|_{2,T}
\le \|(H_2^\flat)'\|_\infty\|e_1\|_{2,T}
\ll\frac{\sqrt{H_1}}{\log T}=o(1).
\end{equation}
In particular, this argument does not apply the finite-polynomial
sampling lemma to an infinite prime-power tail.

For $Q^\flat$, let $\varphi$ be supported on
$[T-3T^{1/2},2T+3T^{1/2}]$ and equal to one on
$[T-2T^{1/2},2T+2T^{1/2}]$.  Both $x^{(1)}_k$ and
$\gamma_k$ lie in this inner interval for sufficiently large $T$:
the former by \cref{lem:multcorrected}, the latter by definition,
and $\Cloc$ is fixed. The segment is therefore contained there.
The pointwise maximal inequality and \cref{lem:maxsample},
now correctly applied with $f=(Q^\flat)'$ and $y_k=x_k^{(1)}$,
give
\begin{equation}\label{eq:Qflattransport}
\|Q^\flat(\gamma)-Q^\flat(x^{(1)})\|_{2,T}
\le2\|e_1\|_{4,T}
\bigl\|\mathcal M(\varphi(Q^\flat)')(x^{(1)})\bigr\|_{4,T}
\ll\frac{\sqrt{H_1}}{\log T}(1+\log X_2)=O(1).
\end{equation}
The coefficient variance used here is
$\sum_{p\le X_2}(\log p)^2/p\ll(1+\log X_2)^2$;
the looser bound $H_0(\log X_2)^2$ is insufficient.
Combining \eqref{eq:tailtransport} and \eqref{eq:Qflattransport}
controls $Q(\gamma)-Q(x^{(1)})$ in $L^2$ by $O(1)$.
Combining the two contributions proves the lemma.
\end{proof}


\section{Dyadic reduction to a deterministic correlation}\label{sec:reduction}
Fix $2\le Y\le X$ and write
\begin{equation}\label{eq:W}
W_{X,Y}(t)=S_X(t)-S_Y(t),\qquad
V_k=W_{X,Y}(\tau_k),\quad U_k=u_k(Y),\quad q_k=u_k(X),\quad h_k=\mathfrak t_k(Y).
\end{equation}
\begin{lemma}[Transport splitting]\label{lem:splitting}
For every $k\in\calI$ and every pair of finite cutoffs,
\begin{equation}\label{eq:transport_split}
\mathfrak t_k(X)=U_k-V_k+(h_k-q_k).
\end{equation}
\end{lemma}
\begin{proof}
Subtract $S_Y$ from $S_X$ in the definition of $\mathfrak t(X)$, and use
$W_{X,Y}(\gamma_k)=S_X(\gamma_k)-S_Y(\gamma_k)=u_k(Y)-u_k(X)$. No interchange of infinite sums is needed: both cutoff sets are finite, and the prime-power series for each fixed prime converges absolutely.
\end{proof}

\begin{theorem}[Quantitative deterministic transfer]\label{thm:transfer}
Retain the finite dyadic window \eqref{eq:windows}. Let $L\ge0$ and suppose that for some numbers $E_U,E_V,E_C\ge0$ and $Q\ge0$,
\begin{equation}\label{eq:transfer_assumptions}
\big|\av{U^2}-L\big|\le E_U,\qquad
\big|\av{V^2}-L\big|\le E_V,\qquad
\big|\av{UV}\big|\le E_C,\qquad
\norm{h-q}\le Q.
\end{equation}
Then
\begin{equation}\label{eq:transfer_bound}
\left|\av{\mathfrak t(X)^2}-2L\right|
\le E_U+E_V+2E_C
+2Q\bigl(\sqrt{L+E_U}+\sqrt{L+E_V}\bigr)+Q^2.
\end{equation}
This statement is uniform in moving cutoffs whenever the four inequalities in \eqref{eq:transfer_assumptions} hold uniformly.
\end{theorem}
\begin{proof}
By \cref{lem:splitting}, the transport equals $U-V+h-q$. Expand its square and average. The contribution from $(U-V)^2$ is $\av{U^2}+\av{V^2}-2\av{UV}$. Its difference from $2L$ is at most $E_U+E_V+2E_C$. Cauchy--Schwarz and the triangle inequality give $\norm{U-V}\le\sqrt{L+E_U}+\sqrt{L+E_V}$, so the remaining linear cross term is bounded by twice this quantity times $Q$; the final quadratic term is bounded by $Q^2$. The proof uses finite sums only.
\end{proof}



\begin{remark}[An endpoint without silent extrapolation]
The algebra of \cref{thm:transfer} does not distinguish $c=1/2$ from smaller fixed positive $c$. The restriction to $c\le1/2$ in a claim based on Dirichlet-polynomial moments must, however, be justified by the degree and total length of the products actually used. The endpoint cannot be obtained by writing ``continuity in $c$'' if the underlying bounds are not uniform there.
\end{remark}


\section{Moving prime bands and exact correlation identities}\label{sec:bands}
For a finite set $\mathcal J$ of primes write
\begin{equation}\label{eq:banddef}
W_{\mathcal J}(t)=\sum_{p\in\mathcal J}a_p\sin(t\log p),
\quad \mu(\mathcal J)=\sum_{p\in\mathcal J}\frac1p,
\quad \mathcal R_p^\tau=\frac2{a_p}\av{s_k\sin(\tau_k\log p)}.
\end{equation}
\begin{proposition}[Band duality, with no limiting process]\label{prop:banddual}
For every finite $\mathcal J$,
\begin{equation}\label{eq:banddual}
\sum_{p\in\mathcal J}\frac{\mathcal R_p^\tau}{p}
=2\pi^2\av{s_kW_{\mathcal J}(\tau_k)}.
\end{equation}
For an arbitrary sequence $b_k$, define $\mathcal B_p(b)=(2/a_p)\av{b_k\sin(\tau_k\log p)}$. Then
\begin{equation}\label{eq:bandb}
\sum_{p\in\mathcal J}\frac{\mathcal B_p(b)}p
=2\pi^2\av{b_kW_{\mathcal J}(\tau_k)}.
\end{equation}
\end{proposition}
\begin{proof}
From $a_p^2=1/(\pi^2p)$ we have $\mathcal R_p^\tau/p=2\pi^2a_p\av{s_k\sin(\tau_k\log p)}$. Sum the finite family. The second identity is identical with $s$ replaced by $b$.
\end{proof}

\begin{theorem}[Uniformity in arbitrary prime subsets]\label{thm:subsets}
Let $\mathscr P_T$ be any finite set of primes, and let $b=(b_k)$. Suppose for every subset $\mathcal J\subseteq\mathscr P_T$ that
\begin{equation}\label{eq:model_hyp}
\left|\av{b_kW_{\mathcal J}(\tau_k)}-
\frac1{2\pi^2}\sum_{p\in\mathcal J}\frac{F_T(p)}p\right|
\le \eta_T\bigl(1+\sqrt{\mu(\mathcal J)}\bigr)
\end{equation}
with a single $\eta_T$ independent of $\mathcal J$. Then
\begin{equation}\label{eq:model_out}
\left|\sum_{p\in\mathcal J}\frac{\mathcal B_p(b)-F_T(p)}p\right|
\le2\pi^2\eta_T\bigl(1+\sqrt{\mu(\mathcal J)}\bigr)
\end{equation}
uniformly for every such subset. If $\mu(\mathscr P_T)\le C$ with fixed $C$, then
\begin{equation}\label{eq:absbound}
\sum_{p\in\mathscr P_T}\frac{|\mathcal B_p(b)-F_T(p)|}p
\le 4\pi^2\eta_T(1+\sqrt C).
\end{equation}
\end{theorem}
\begin{proof}
Apply \cref{prop:banddual} to \eqref{eq:model_hyp}. For the absolute-value estimate, partition $\mathscr P_T$ into the subset of primes with nonnegative real difference $\mathcal B_p(b)-F_T(p)$ and its complement. All quantities are real. The sum of absolute values is the first signed sum minus the second, and each is bounded by \eqref{eq:model_out}. This argument requires the \emph{same} uniform estimate for all subsets, including subsets selected by the signs; validity only for consecutive intervals does not suffice.
\end{proof}


\section{A finite independent-prime comparison}\label{sec:model}
Let $\Theta_p$ be independent random variables uniformly distributed on $[0,2\pi]$. For $p\le Y$ define
\begin{equation}\label{eq:phase_model}
f_p(\vartheta)=-\frac1\pi\sum_{m\ge1}\frac{\sin(m\vartheta)}{m p^{m/2}},
\quad P_Y=\sum_{p\le Y}f_p(\Theta_p),
\quad \Phi_Y(\xi)=\mathbb E\e^{-i\xi P_Y}.
\end{equation}
\begin{lemma}[Characteristic factor and its variance]\label{lem:char}
For $|\xi|\le1$,
\begin{equation}\label{eq:char_product}
\Phi_Y(\xi)=\prod_{p\le Y}\frac1{2\pi}\int_0^{2\pi}\e^{-i\xi f_p(\vartheta)}\,d\vartheta,
\end{equation}
and, uniformly as $Y\to\infty$,
\begin{equation}\label{eq:char_asymp}
\log\Phi_Y(\xi)=-\frac{\xi^2}{4\pi^2}\log\log Y+O(\xi^2+\xi^4).
\end{equation}
For sufficiently small $|\xi|$ the logarithm is the real logarithm. The $O(\xi^2)$ includes a convergent arithmetic constant and must not be discarded in a fixed-frequency asymptotic. This lemma concerns complete Euler factors, whereas the sampled profile later in the paper contains truncated prime-power polynomials; their positivity and fourth-order expansion are verified separately.
\end{lemma}
\begin{proof}
Independence proves \eqref{eq:char_product}. Since $f_p$ is odd, each factor equals the average of $\cos(\xi f_p)$. Orthogonality of sine harmonics gives
\begin{equation}\label{eq:variance}
\mathbb E f_p(\Theta_p)^2=\frac1{2\pi^2}\sum_{m\ge1}\frac1{m^2p^m}=
\frac1{2\pi^2p}+O(p^{-2}).
\end{equation}
Moreover $\|f_p\|_\infty\le\pi^{-1}\sum_{m\ge1}(mp^{m/2})^{-1}\ll p^{-1/2}$, hence
$\mathbb E\cos(\xi f_p)=1-\xi^2\mathbb Ef_p^2/2+O(\xi^4p^{-2})$. Each factor is strictly positive for $|\xi|\le1$ because $|\xi f_p|<\pi/2$ (in fact $\sup|f_p|<1$). Taking logarithms, summing, and using $\sum_{p\le Y}p^{-1}=\log\log Y+B_1+o(1)$ gives the assertion.
\end{proof}

\begin{remark}[Model versus deterministic clock]
\cref{lem:char} concerns independent phases $\Theta_p$. To replace its right side by a clock average of prime polynomials evaluated at $\tau_k$ requires a separate bound on all nonzero logarithmic frequencies appearing after expansion. The length of the multiplied Dirichlet polynomials, and not merely the largest individual prime cutoff, determines whether mean-value estimates apply. This remains true when $\xi$ depends on a moving prime $p=T^\theta$.
\end{remark}


\section{A deterministic two-scale phase expansion}\label{sec:twoscale}
For fixed $0<c\le1/2$ set
\begin{equation}\label{eq:transparams}
\begin{gathered}
V_0=\frac{H_0}{2\pi^2},\quad
\omega=\frac1\pi\sqrt{\log\log V_0},\quad
\theta_T=\frac{\omega}{\sqrt{V_0}},\\
Y=T^{\theta_T},\quad Z=T^{H_0^{-2}},\quad X=T^c,\\
W=S_X-S_Y,\qquad Q=S_{X_2}.
\end{gathered}
\end{equation}
For sufficiently large $T$ (depending on $c$),
$Z<X_1<Y<X_2<X\le T^{1/2}$.  Define
\begin{equation}\label{eq:bandenergy}
L=\frac1{2\pi^2}\log\frac c{\theta_T}
 =\frac{H_1}{4\pi^2}+O_c(1+\log H_2).
\end{equation}
These parameter relations are asymptotic: no assertion concerning
small or moderate heights is made by invoking them.

\begin{proposition}[Reduction of the phase energy to one correlation]
\label{prop:reduceA}
Put
\[
 B_c(T)=\av{[S_X(\gamma_k)-S_X(\tau_k)]^2}.
\]
Then
\begin{equation}\label{eq:reduceA}
B_c(T)=2L-2\av{s_kW(\tau_k)}
 +O_c(\sqrt{H_1H_2}).
\end{equation}
\end{proposition}
\begin{proof}
The identity $W(\gamma)=u(Y)-u(X)$ and
$\mathfrak t(X)=\mathfrak t(Y)+W(\gamma)-W(\tau)$ is exact.
By \cref{thm:primedefect},
$\|u(Y)\|_2^2=L+O(1+\sqrt{H_1})$ and
$\|u(X)\|_2=O_c(1)$.  Consequently
$\av{W(\gamma)^2}=L+O_c(\sqrt{H_1})$.
The almost-orthogonality theorem and the negligible prime-power
part give $\av{W(\tau)^2}=L+o(1)$.
\cref{lem:shorttransport} gives
$\|\mathfrak t(Y)\|_2\ll\sqrt{H_0}\theta_T
\ll\omega\ll\sqrt{H_2}$.  Hence the cross term involving this
short transport costs $O(\omega\sqrt L+\omega^2)$.
Finally,
$W(\gamma)=s-u(X)-S_Y(\gamma)$ and
$S_Y(\gamma)=S_Y(\tau)+\mathfrak t(Y)$.
Thus
\[
\av{W(\gamma)W(\tau)}
 =\av{sW(\tau)}+O_c(\sqrt L+\omega\sqrt L),
\]
using Cauchy--Schwarz, the small norm of $u(X)$, and the
clock orthogonality of the disjoint prime bands of $S_Y$ and $W$.
The remaining prime powers in $S_Y$ have uniformly bounded
mean-square norm, giving $O(\sqrt L)$ as claimed.
Expanding the square proves \eqref{eq:reduceA}.
\end{proof}

\begin{proposition}[Removing the zeros from the correlation]
\label{prop:reduceB}
\begin{equation}\label{eq:reduceB}
\av{s_kW(\tau_k)}=
\av{Q(x_k^{(1)})W(\tau_k)}
 +O(\sqrt{H_1(1+H_2)}).
\end{equation}
\end{proposition}
\begin{proof}
\cref{lem:secondreplacement} bounds the $L^2$ norm of
$s-Q(x^{(1)})$ by $O(1+\sqrt{H_2})$.
The prime-band orthogonality and the prime-power estimate give
$\|W(\tau)\|_2\ll\sqrt L\ll\sqrt{H_1}$.
Cauchy--Schwarz yields the stated estimate.  In particular the
right-hand side is an arithmetic correlation evaluated entirely at
points computable from the smooth clock and primes.
\end{proof}


\subsection{Finite prime-power reduction}
Write $\beta_{p^m}=-1/(\pi m p^{m/2})$.  Introduce the
finite polynomials
\begin{equation}\label{eq:truncations}
\begin{split}
Q^\flat(t)&=\sum_{p\le X_2}\sum_{p^m\le X_2^2}
 \beta_{p^m}\sin(mt\log p),\\
P_0(t)&=\sum_{p\le Z}\sum_{p^m\le Z^2}
 \beta_{p^m}\sin(mt\log p),\\
P_1(t)&=\sum_{Z<p\le X_1}\sum_{p^m\le X_1^2}
 \beta_{p^m}\sin(mt\log p),\\
P_2(t)&=\sum_{p\le Z}\sum_{Z^2<p^m\le X_1^2}
 \beta_{p^m}\sin(mt\log p),\\
E_1(t)&=S_{X_1}(t)-(P_0+P_1+P_2)(t),\\
W_{\rm pr}(t)&=\sum_{Y<p\le X}a_p\sin(t\log p),
\qquad W_{\rm pp}(t)=W(t)-W_{\rm pr}(t).
\end{split}
\end{equation}
The distinction between primes $p\le X$ and prime powers
$p^m\le X^2$ in these truncations is deliberate.

\begin{lemma}[Reduction without an unstated limit exchange]
\label{lem:finitecorr}
Let
$y_k=\tau_k-[P_0(\tau_k)+P_1(\tau_k)]/\Lambda(\tau_k)$.
Then, uniformly for fixed $0<c\le1/2$,
\begin{equation}\label{eq:finitecorr}
\av{Q(x_k^{(1)})W(\tau_k)}
=\av{Q^\flat(y_k)W_{\rm pr}(\tau_k)}+o(1).
\end{equation}
\end{lemma}
\begin{proof}
For $p\le R$ the absolute sum of the harmonics with
$p^m>R^2$ is $O(R^{-1})$ for each prime, and summing those tails
using the reciprocal-prime estimate gives
\begin{equation}\label{eq:finiteerrs}
\|Q-Q^\flat\|_\infty\ll1/\log X_2,
\qquad\|E_1\|_\infty\ll1/\log X_1.
\end{equation}
Since $\|W(\tau)\|_2\ll\sqrt{H_1}$, replacing
$Q$ by $Q^\flat$ costs $O(\sqrt{H_1}/\log X_2)=o(1)$.

For the prime powers in $W$, the Sobolev averaging inequality
on the disjoint clock cells and Montgomery--Vaughan give
$\|W_{\rm pp}(\tau)\|_2\ll Y^{-1/4}$.
By \cref{lem:multcorrected,lem:polysampling} applied to
$Q^\flat$, with $q=1$ and $R=X_2^2=T^{2/\sqrt{H_1}}$,
$\|Q^\flat(x^{(1)})\|_2\ll\sqrt{H_0}$.
Their product is $O(\sqrt{H_0}Y^{-1/4})=o(1)$.
All length conditions hold for $T$ sufficiently large because
$R^{8q}(2\|Q^\flat\|_1)^{4q}=T^{o(1)}$.

For the last replacement put
$\Pi_v=P_0+P_1+v(P_2+E_1)$ and
$y_k(v)=\tau_k-\Pi_v(\tau_k)/\Lambda(\tau_k)$.
The fundamental theorem of calculus gives the exact formula
\[
Q^\flat(x^{(1)}_k)-Q^\flat(y_k)
=-\frac{(P_2+E_1)(\tau_k)}{\Lambda(\tau_k)}
\int_0^1(Q^\flat)'(y_k(v))dv.
\]
The coefficient variance of $P_2$ is $O(Z^{-1})$:
each individual prime contributes at most $O(Z^{-1})$ in
absolute coefficient sum, and there are $O(Z)$ such primes.
The clock-model comparison of \cref{lem:klclock} and the
independent-prime fourth moment give
$\|P_2(\tau)\|_4\ll Z^{-1/2}+T^{-1/8}$.
Uniformly in $v$, \cref{lem:multcorrected,lem:polysampling}
give
$\|(Q^\flat)'(y(v))\|_4\ll\sqrt{H_0}\log X_2$.
By Holder with exponents $(4,4,2)$, the resulting correlation
error is
\[
\ll \frac{\sqrt{H_0}\log X_2}{\log T}
\left(Z^{-1/2}+\frac1{\log X_1}\right)\sqrt{H_1}=o(1).
\]
The estimates are uniform in $v\in[0,1]$; integrating in $v$
therefore introduces no unverified interchange.
\end{proof}

\subsection{Two-scale expansion with integral remainders}
For integers $J\ge1$ write
$e_J(z)=\sum_{j=0}^{J-1}z^j/j!$ and
$R_J(z)=e^z-e_J(z)$.
The following formula is used only with the finite polynomial
$Q^\flat$, so it involves no infinite series expansion.

\begin{lemma}[Exact aggregated Taylor formula]\label{lem:aggregated}
Let $P_i=P_i(\tau_k)$ and $\Lambda=\Lambda(\tau_k)$.
For every pair of positive integers $J_0,J_1$,
\begin{equation}\label{eq:aggregated}
Q^\flat(y_k)=\mathsf{Main}_k+
\mathsf{Rem}_{I,k}+\mathsf{Rem}_{II,k},
\end{equation}
where
\begin{align}
\mathsf{Main}_k
 &=\sum_{j<J_0}\sum_{\ell<J_1}
 \frac{(-P_0)^j(-P_1)^\ell}{j!\ell!\Lambda^{j+\ell}}
 (Q^\flat)^{(j+\ell)}(\tau_k),\label{eq:aggmain}\\
\mathsf{Rem}_{I,k}
 &=\frac{(-P_1)^{J_1}}{(J_1-1)!\Lambda^{J_1}}
 \int_0^1(1-v)^{J_1-1}(Q^\flat)^{(J_1)}
 \left(\tau_k-\frac{P_0+vP_1}{\Lambda}\right)dv,\label{eq:aggI}\\
\mathsf{Rem}_{II,k}
 &=\sum_{\ell<J_1}
 \frac{(-P_1)^\ell(-P_0)^{J_0}}
 {\ell!(J_0-1)!\Lambda^{\ell+J_0}}
 \int_0^1(1-v)^{J_0-1}(Q^\flat)^{(\ell+J_0)}
 \left(\tau_k-\frac{vP_0}{\Lambda}\right)dv.\label{eq:aggII}
\end{align}
\end{lemma}
\begin{proof}
Each summand of $Q^\flat$ is a sine harmonic of frequency
$\log n$.  Its value at $y_k$ is
$\Im(n^{i\tau_k}e^{a+b})$, where
$a=-i(\log n)P_0/\Lambda$ and
$b=-i(\log n)P_1/\Lambda$.
The elementary identity
\[
e^{a+b}=e_{J_0}(a)e_{J_1}(b)
 +e^aR_{J_1}(b)+R_{J_0}(a)e_{J_1}(b)
\]
is exact.  The integral remainder formula
$R_J(z)=z^J\int_0^1(1-v)^{J-1}e^{vz}dv/(J-1)!$
turns the last two terms into derivatives of the entire
polynomial $Q^\flat$ evaluated at the displayed shifted points.
Summing the finitely many harmonics produces
\eqref{eq:aggregated}--\eqref{eq:aggII}.  In particular the
untruncated factor $e^a$ is retained inside a shifted evaluation;
it is not bounded by its absolute coefficient sum.
\end{proof}

\begin{lemma}[Uniform control of both remainders]\label{lem:aggerrors}
Let
\begin{equation}\label{eq:aggparameters}
\widehat\xi=\frac{\log X_2^2}{\Lambda(T-\Cloc-1)},
\quad \sigma_0^2=\sum_{p\le Z}
\left(\sum_{p^m\le Z^2}|\beta_{p^m}|\right)^2,
\quad \sigma_1^2=\sum_{Z<p\le X_1}
\left(\sum_{p^m\le X_1^2}|\beta_{p^m}|\right)^2,
\end{equation}
and choose
\[
J_1=\lceil H_1\rceil,\qquad
J_0=\lceil32e^2\widehat\xi^2\sigma_0^2\rceil.
\]
For these choices,
\begin{equation}\label{eq:aggerrors}
\av{\left(|\mathsf{Rem}_I|+|\mathsf{Rem}_{II}|\right)
 |W_{\rm pr}(\tau)|}=o(1).
\end{equation}
The estimate is uniform when the terminal prime interval is replaced
by an arbitrary subset of $(X_1,T^{1/2}]$ and
$\sqrt L$ is replaced by its harmonic square-root mass.
\end{lemma}
\begin{proof}
The cutoffs give
$\widehat\xi=(4\pi/\sqrt{H_1})(1+o(1))$,
$\sigma_0^2\le(H_0/\pi^2)(1+o(1))$ and
$\sigma_1^2\le(3H_1/(2\pi^2))+O(1)$.
Thus $\widehat\xi^2\sigma_1^2\le25$ and
$J_0\asymp H_0/H_1$.
For $P_1$ and $P_0$ evaluated at $\tau_k$, all even moments
needed below agree with the independent-prime model up to
$T^{-1/2}$: expand their powers as rational trigonometric
polynomials, use \cref{lem:klclock}, and check
\[
(X_1^2)^{4J_1}=T^{8H_1/\sqrt{H_0}}=T^{o(1)},
\qquad (Z^2)^{8J_0}=T^{16J_0/H_0^2}=T^{o(1)}.
\]
Hoeffding's moment bound in the model reads
$\mathbb E|P_i|^{2q}\le2q!(2\sigma_i^2)^q$.

For \eqref{eq:aggI} normalize
$\widetilde Q^{(j)}=(Q^\flat)^{(j)}/(\log X_2^2)^j$.
Its coefficients are dominated by those of $Q^\flat$,
independently of $j$.  Jensen's inequality for the probability
measure $J_1(1-v)^{J_1-1}dv$, Holder with exponents
$(4,4,2)$ and \cref{lem:multcorrected,lem:polysampling}
give
\[
\av{|\mathsf{Rem}_I W_{\rm pr}|}
\ll \left\|\frac{(\widehat\xi|P_1|)^{J_1}}{J_1!}\right\|_4
\sup_{v}\|\widetilde Q^{(J_1)}(y(v))\|_4
\|W_{\rm pr}(\tau)\|_2.
\]
The last two factors are respectively $O(\sqrt{H_0})$
and $O(\sqrt L)$.  The model moment bound and the factorial
inequality $J!\ge(J/e)^J$ give
\[
\left\|\frac{(\widehat\xi P_1)^{J_1}}{J_1!}\right\|_4
 \le2\left(\frac{100e^2}{J_1}\right)^{J_1/2}.
\]
As $J_1\sim H_1=\log H_0$, the first remainder contributes
$O(H_0^{1/2-(\log H_1)/2+O(1)}\sqrt L)=o(1)$.

For \eqref{eq:aggII}, use Holder with exponents
$(8,8,4,2)$ and sum over $\ell<J_1$.  The model moment bound
implies
\[
\sum_{\ell<J_1}\left\|\frac{(\widehat\xi P_1)^\ell}{\ell!}\right\|_8
 =O(1),
\qquad
\left\|\frac{(\widehat\xi P_0)^{J_0}}{J_0!}\right\|_8
 \le2^{1-J_0}.
\]
For the second estimate use $J_0\ge32e^2\widehat\xi^2\sigma_0^2$
and the eighth-moment factorial estimate; every required model
moment is within the frequency-height range checked above.
By bounded multiplicity and the sampling lemma,
$\|\widetilde Q^{(\ell+J_0)}(\tau-vP_0/\Lambda)\|_4
\ll\sqrt{H_0}$ uniformly in $\ell,v$.
Thus the second remainder is
$O(2^{-J_0}\sqrt{H_0L})=o(1)$.
There is no dependence on whether the primes of the terminal band
are consecutive: only their coefficient-square mass appears.
\end{proof}

\begin{lemma}[All nonzero rational frequencies]\label{lem:aggoffdiag}
Let $\mathsf{Main}$ be given by \eqref{eq:aggmain}.
With $0<c\le1/2$ fixed, the difference between its average
against $W_{\rm pr}(\tau_k)$ and its independent-prime diagonal is
\begin{equation}\label{eq:aggoffdiag}
O\left(T^{3c/2-1+o(1)}\right)=o(1).
\end{equation}
The estimate is uniform in all subsets of the terminal prime band,
including when $c=1/2$.
\end{lemma}
\begin{proof}
Expand the finitely many $P_0,P_1,Q^\flat,W_{\rm pr}$
factors into exponentials.  Each frequency is a rational
logarithmic frequency $\log(a/b)$, where both numerator and
denominator are bounded by the product of the individual
lengths.  At Taylor orders $j<J_0$, $\ell<J_1$ this height is
\[
\le Z^{2J_0}X_1^{2J_1}X_2^2X
 =T^{c+\eta_T},\qquad
\eta_T=\frac{2J_0}{H_0^2}+
\frac{2J_1}{\sqrt{H_0}}+\frac2{\sqrt{H_1}}=o(1).
\]
The weight $\Lambda(\tau_k)^{-(j+\ell)}$ is monotone
and has supremum plus variation at most twice its value at
$T-1$.  Applying \cref{lem:klclock} after partial summation
to each nonzero frequency, and dividing by $M_T$, costs at most
$O(T^{c+\eta_T-1})$ times its coefficient absolute value.
The absolute coefficient sum of the entire expansion is at most
\[
J_0J_1(2Z\widehat\xi)^{J_0}
 (2X_1\widehat\xi)^{J_1}
 \cdot 2X_2\sqrt X=T^{c/2+o(1)}.
\]
Hence the total nonzero-frequency contribution is
$T^{3c/2-1+o(1)}$.  For $c=1/2$ this is
$T^{-1/4+o(1)}$, strictly decaying.  The calculation is unchanged
for any subset, since removing primes can only reduce the
absolute coefficient sum and the maximal height.
\end{proof}

\begin{lemma}[The diagonal from unique factorization]\label{lem:aggdiag}
Let
$\xi_{p,k}=\log p/\Lambda(\tau_k)$.
Then the diagonal of $\mathsf{Main}\,W_{\rm pr}$ is exactly
\begin{equation}\label{eq:aggdiag}
\sum_{Y<p\le X_2}\frac{a_p^2}{2}
\av{\Re\,\mathbb E_{\rm mod}
 [e_{J_0}(-i\xi_{p,k}P_0)]
 \mathbb E_{\rm mod}[e_{J_1}(-i\xi_{p,k}P_1)]}.
\end{equation}
It is bounded by
\begin{equation}\label{eq:diagdecay}
O\left(L\exp[-2\pi^2\omega^2(1+o(1))]
 +L2^{-J_0}\right)=o(1).
\end{equation}
\end{lemma}
\begin{proof}
Use independent uniform phases $\Theta_p$ and the associated
prime-power harmonics $e^{im\Theta_p}$.
The prime factors of $P_0,P_1$ are $\le X_1<Y$, whereas those of
$W_{\rm pr}$ lie above $Y$.  A product has constant total frequency
only when the prime $p$ selected from $W_{\rm pr}$ appears in
$Q^\flat$ with exponent one.  Thus $p\le X_2$, and its sine
inner product contributes exactly $a_p^2/2$.
All other prime phases cancel independently by unique factorization.
The sum over the Taylor orders gives \eqref{eq:aggdiag}.

For imaginary arguments the remainder satisfies
$|R_J(iz)|\le|z|^J/J!$.  The moment bounds used in
\cref{lem:aggerrors} show that the truncated expectations
may be replaced by the full characteristic functions at a cost
$O(2^{-J_0}+(50e/J_1)^{J_1/2})$, uniformly in $k,p$.
For $P_0$, independence across primes and the second-moment
calculation
$\mathbb E f_r(\Theta_r)^2=(2\pi^2r)^{-1}+O(r^{-2})$
give, when $p\ge Y$,
\[
\left|\mathbb E_{\rm mod}e^{-i\xi_{p,k}P_0}\right|
\le\exp\left[-\frac{\xi_{p,k}^2}{4\pi^2}
  (H_0-2H_1+O(1))+O(\xi_{p,k}^3)\right].
\]
To justify the logarithm factorwise, use the uniform bound
\eqref{eq:truncatedphasebound}, valid for the truncated prime-power
phases as well as for the full phase.  Since the relevant $\xi$
satisfies $0\le\xi\le\pi+o(1)$, the individual cosine factors
are strictly positive; their quadratic Taylor expansions have a
summable fourth-order error $O(\xi^4r^{-2})$.  In the present
diagonal, the stronger condition $p\le X_2$ even gives
$\xi\ll H_1^{-1/2}=o(1)$.  Now
$\xi_{p,k}\ge(2\pi\theta_T)(1+o(1))$.
Using $\theta_T^2H_0=2\pi^2\omega^2$ and
$\theta_T^2H_1\to0$, the displayed exponential is at most
$\exp[-2\pi^2\omega^2(1+o(1))]$.
The $P_1$ characteristic factor has modulus at most one.
The total prime diagonal mass is at most $L+o(1)$.
Because $\omega^2=\pi^{-2}\log\log V_0$,
$L\exp[-2\pi^2\omega^2(1+o(1))]
\ll H_1(\log V_0)^{-2+o(1)}=o(1)$.
This proves \eqref{eq:diagdecay}.
\end{proof}

\begin{theorem}[Deterministic clock correlation]\label{thm:clockcorr}
For $0<c\le1/2$ fixed,
\begin{equation}\label{eq:clockcorr}
\av{S_{X_2}(x_k^{(1)})
 [S_{T^c}(\tau_k)-S_{T^{\theta_T}}(\tau_k)]}=o(1).
\end{equation}
\end{theorem}
\begin{proof}
Apply \cref{lem:finitecorr}; expand its right side by
\cref{lem:aggregated}; bound the two integral remainders by
\cref{lem:aggerrors}, all nonzero frequencies by
\cref{lem:aggoffdiag}, and the unique-factorization diagonal
by \cref{lem:aggdiag}.  Every interchange of sums is over
finitely supported prime-power polynomials; the only integral
interchanges are over $v\in[0,1]$, where the uniform majorants
established in \cref{lem:aggerrors} apply.  The result follows.
\end{proof}



\section{Transport energy and varying prime cutoffs}\label{sec:energy}
The estimates of \cref{sec:twoscale} yield the fixed-cutoff transport
energy and its dependence on the moving cutoff.  The distinction
between a cutoff fixed as an exponent and a cutoff fixed as an integer
is essential.

\begin{theorem}[Transport energy at a fixed power cutoff]\label{thm:energy}
For every fixed $c\in(0,\tfrac12]$, set $X=T^c$ and let
$\mathcal I_T$ be the true-ordinate window \eqref{eq:windows}.
Then, as $T\to\infty$,
\begin{equation}\label{eq:energy}
\av{\bigl[S_{T^c}(\gamma_k)-S_{T^c}(\tau_k)\bigr]^2}
=\frac{\log\log\log T}{2\pi^2}
+O_c\!\left(\sqrt{\log\log\log T\,
                        \log\log\log\log T}\right),
\end{equation}
and the same equivalent, with the same error order, holds for
$\av{r_k(T^c)^2}$.
\end{theorem}
\begin{proof}
By \cref{prop:reduceA,prop:reduceB,thm:clockcorr},
\[
B_c(T)=2L+O_c(\sqrt{H_1H_2})
\]
(the product of the $L^2$ replacement error with the band
norm is already included in \cref{prop:reduceB}).
The exact choice of scale in \eqref{eq:transparams} gives
\[
2L=\frac1{\pi^2}\log\frac c{\theta_T}
=\frac{H_1}{2\pi^2}+O_c(1+\log H_2).
\]
The $O_c(1+\log H_2)$ error is absorbed in
$O_c(\sqrt{H_1H_2})$ for all sufficiently large $T$.
This proves \eqref{eq:energy}.

For the second assertion use the exact displacement decomposition
\eqref{eq:exact_r}.  By \cref{thm:primedefect},
$\|u(T^c)\|_2=O_c(1)$, and by Taylor's theorem
$\|e\|_2\ll T^{-1}$ because the bare-clock displacement is bounded.
The reverse triangle inequality gives
$|\|r\|_2-\|\mathfrak t\|_2|=O_c(1)$.
Squaring and using $\|\mathfrak t\|_2\asymp\sqrt{H_1}$ adds
at most $O_c(\sqrt{H_1})$, smaller than the displayed error.
\end{proof}

\begin{corollary}[Moving exponents close on the logarithmic scale]
\label{cor:varenergy}
If $c=c(T)$ satisfies $0<c\le1/2$ and
\begin{equation}\label{eq:varcondition}
\log(1/c)=o(H_1),
\end{equation}
then
\begin{equation}\label{eq:varenergy}
\av{\mathfrak t_k(T^c)^2}
\sim\av{r_k(T^c)^2}
\sim\frac{H_1}{2\pi^2}.
\end{equation}
No equivalent is asserted outside \eqref{eq:varcondition}.
\end{corollary}
\begin{proof}
The hypothesis implies $c>\theta_T$ eventually.
By \cref{thm:primedefect}, $\|u(T^c)\|_2^2=o(H_1)$.
The short-band phase $u(Y)$ has square norm
$H_1/(4\pi^2)+o(H_1)$, and therefore
$\|W(\gamma)\|_2^2=H_1/(4\pi^2)+o(H_1)$.
The clock-band square norm has the same equivalent by
\cref{cor:massbands} and the prime-power estimate.
The deterministic clock correlation estimate
\cref{thm:clockcorr} is uniform for $\theta_T\le c\le1/2$:
in \cref{lem:aggoffdiag} the saving is bounded by
$T^{-1/4+o(1)}$, and all remainder bounds depend only on the
harmonic mass, at most $O(H_1)$.
The cross terms involving $u(T^c)$ are $o(H_1)$ by
Cauchy--Schwarz.  The short transport costs $o(H_1)$ as before.
Thus $B_c=H_1/(2\pi^2)+o(H_1)$.
Finally \eqref{eq:exact_r} and $\|u(T^c)\|_2=o(\sqrt{H_1})$
give the same equivalent for $r$.
\end{proof}

\begin{corollary}[Mean-square distance in ordinate height]
\label{cor:heightdistance}
For fixed $c\in(0,1/2]$,
\begin{equation}\label{eq:heightdistance}
\av{\bigl(\tau_k^{T^c}-\gamma_k\bigr)^2}
=\frac{2H_1+O_c(\sqrt{H_1H_2})}{\log^2T}.
\end{equation}
\end{corollary}
\begin{proof}
By definition
$\tau_k^{T^c}-\gamma_k=r_k(T^c)/\Lambda(\tau_k)$.
The unconditional localization \cref{lem:clocklocal} places $\tau_k\in[T-O(1),2T+O(1)]$,
so $\Lambda(\tau_k)^{-2}
=(2\pi/\log T)^2(1+O(1/\log T))$ uniformly
in the true-ordinate window.  Multiply the second estimate of
\cref{thm:energy} by this factor; the resulting constant is
$4\pi^2/(2\pi^2)=2$.
\end{proof}

\begin{proposition}[A lower bound for intermediate cutoffs]
\label{prop:intermediate}
Uniformly for $\theta_T\le c\le1/2$,
\begin{equation}\label{eq:lowerenergy}
\av{r_k(T^c)^2}
\ge\frac1{2\pi^2}\log\frac c{\theta_T}
 -O(\sqrt{H_1H_2}).
\end{equation}
For every fixed $\alpha$ with $0<\alpha<1/2$, the choice
$c=H_0^{-\alpha}$ satisfies the required range for sufficiently
large $T$ and yields
\begin{equation}\label{eq:alphalower}
\av{r_k(T^{H_0^{-\alpha}})^2}
\ge\frac{\frac12-\alpha}{2\pi^2}H_1-o(H_1),
\qquad\av{r_k(T^{H_0^{-\alpha}})^2}\ll_\alpha H_1.
\end{equation}
\end{proposition}
\begin{proof}
Let $W_{\rm pr}=\sum_{Y<p\le T^c}a_p\sin(t\log p)$
and $L=(2\pi^2)^{-1}\log(c/\theta_T)\ge0$.
If $L\ll\sqrt{H_1H_2}$, \eqref{eq:lowerenergy} is trivial
because the left side is nonnegative.  Otherwise
$\|W_{\rm pr}(\tau)\|_2^2=L+o(1)$.
By the replacement argument of \cref{prop:reduceB} and the
uniform diagonal of \cref{lem:aggdiag},
$\av{sW_{\rm pr}(\tau)}=O(\sqrt{H_1H_2})$.
Orthogonality of disjoint clock bands and the $L^2$ boundedness
of the prime-power part give
$\av{S_{T^c}(\tau)W_{\rm pr}(\tau)}
 =L+O(\sqrt L)+o(1)$.
The exact identity $r=s-S_{T^c}(\tau)+O(T^{-1})$ therefore gives
$\av{rW_{\rm pr}(\tau)}=-L+O(\sqrt{H_1H_2})$.
Cauchy--Schwarz gives
\[
\av{r^2}\ge
\frac{|\av{rW_{\rm pr}}|^2}{\av{W_{\rm pr}^2}}
\ge L-O(\sqrt{H_1H_2}),
\]
where the final estimate is applied only in the nontrivial range
of $L$.  For $c=H_0^{-\alpha}$, the decomposition
$r=u(T^c)+\mathfrak t(Y)+W(\gamma)-W(\tau)+O(T^{-1})$
and the established second moments give $\av{r^2}\ll_\alpha H_1$.
The case $\alpha=0$ is excluded from \eqref{eq:alphalower}:
it would give $c=1$, outside the proved $c\le1/2$ range.
\end{proof}



\section{Prime-power projection geometry and index-dependent corrections}
\label{sec:variableprojection}
We now extend the fixed-exponent conclusions to a one-step correction whose prime cutoff can change separately for every index. All results in this section use the \emph{clock index window}
\begin{equation}\label{eq:projwindow}
\mathcal J_T=\{k:T\le\tau_k<2T\},\qquad M=|\mathcal J_T|,\qquad
\avgJ{z_k}=\frac1M\sum_{k\in\mathcal J_T}z_k.
\end{equation}
The smooth clock and all trigonometric phases are the same as above. We do not assume that the values $X_k$ or the corrected ordinates are monotone. The point of using $\mathcal J_T$ is that consecutive clock points are separated, permitting a direct maximal argument. The transfer to the true-ordinate window is stated explicitly below.

\subsection{Prime-power phases and operator-norm almost orthogonality}
Here $N_0$ and the midpoint ordinates $\tau_k$ are as in \eqref{eq:N0}. Let $\tau_k$ be the eventually increasing sequence defined by
$N_0(\tau_k)=k-\tfrac12$ (the finitely many initial values do not affect the theorem). Fix $0<b<1$, let
\begin{equation}
\mathcal Q_b(T)=\{p^a: p\text{ prime},\ a\in\mathbb N,\ 2\le p^a\le T^b\},\quad
\mathcal J_T=\{k:T\le\tau_k<2T\},\quad M=|\mathcal J_T|,
\label{pp:eq:sets}
\end{equation}
and introduce vectors in $\mathbb R^M$ by
\begin{equation}
 v_n(k)=\sqrt{\frac{2}{M}}\sin(\tau_k\log n),\qquad
 G=(\langle v_n,v_m\rangle)_{n,m\in\mathcal Q_b(T)}.
\label{pp:eq:gram}
\end{equation}
The clock-index window $\mathcal J_T$ differs from the true-ordinate window $\mathcal I_T$ used earlier; their boundary comparison is proved below. All implied constants may depend on fixed $b$ but not on $T$. The number of points satisfies
\begin{equation}
M=N_0(2T)-N_0(T)+O(1)\asymp T\log T.
\label{pp:eq:M}
\end{equation}

\subsection{Exponential sums on the smooth clock}
Set, for $u>0$,
\begin{equation}
C_T(u)=\frac1M\sum_{k\in\mathcal J_T}e^{i\tau_k\log u}.
\label{pp:eq:C}
\end{equation}
For the following estimates we use two standard unconditional exponential-sum inequalities. If a $C^1$ real phase $F$ has monotone derivative whose distance to the nearest integer is at least $\lambda>0$, the Kusmin--Landau bound is $|\sum e^{2\pi i F(k)}|\ll\lambda^{-1}$ (when the derivative remains in one component of $\mathbb R\setminus\mathbb Z$). If $F''$ has constant sign and $\lambda\le |F''|\le A\lambda$, with fixed $A$, the van der Corput second-derivative bound is
\begin{equation}
\left|\sum_{k\in J\cap\mathbb Z}e^{2\pi i F(k)}\right|
\ll_A |J|\sqrt\lambda+\lambda^{-1/2}.
\label{pp:eq:vdC}
\end{equation}
Both estimates hold for any real interval $J$ of consecutive indices.

Extend $k\mapsto\tau_k$ to a smooth increasing inverse $x\mapsto\tau_x$ of $N_0$, and for $u>1$ put
\begin{equation}
F_u(x)=\frac{\tau_x\log u}{2\pi}.
\label{pp:eq:F}
\end{equation}
Differentiation of $N_0(\tau_x)=x-\tfrac12$ gives, uniformly when $\tau_x\in[T,2T]$,
\begin{equation}
F_u'(x)=\frac{\log u}{\log(\tau_x/2\pi)},\qquad
F_u''(x)=-\frac{2\pi\log u}{\tau_x\log^3(\tau_x/2\pi)}.
\label{pp:eq:derivs}
\end{equation}
Thus $F_u'$ is monotone decreasing and $F_u''<0$ for $u>1$.

\begin{lemma}[Quotient frequencies]\label{pp:lem:quotient}
If $2\le m<n\le T^b$ are integers (not necessarily prime powers), then
\begin{equation}
|C_T(n/m)|\ll_b\frac{1}{T\log(n/m)}
\ll_b\frac{\max(n,m)}{T|n-m|}.
\label{pp:eq:quotient}
\end{equation}
\end{lemma}
\begin{proof}
The derivative in \eqref{pp:eq:derivs} with $u=n/m$ obeys
\[
0<\frac{\log(n/m)}{\log(T/\pi)}\le F_u'(x)
\le \frac{b\log T}{\log(T/2\pi)}<\frac{1+b}{2}<1
\]
for $T$ large depending on $b$. Hence its distance to an integer is at least a constant depending on $b$ times $\log(n/m)/\log T$. Kusmin--Landau and \eqref{pp:eq:M} give the first bound. Finally
$\log(n/m)=\int_m^n dt/t\ge(n-m)/n$.
\end{proof}

\begin{lemma}[Product frequencies]\label{pp:lem:product}
For integers $n,m\ge2$ with $nm\le T^{2b}$, write $u=nm$. Then
\begin{align}
|C_T(u)|&\ll_b T^{-1},&&u\notin[T/(4\pi),4T],\label{pp:eq:nonres}\\
|C_T(u)|&\ll T^{-1/2},&&u\in[T/(4\pi),4T].\label{pp:eq:res}
\end{align}
\end{lemma}
\begin{proof}
For $u<T/(4\pi)$, $0<F_u'(x)<1$ and
\[
1-F_u'(x)
=\frac{\log(\tau_x/(2\pi u))}{\log(\tau_x/2\pi)}
\ge\frac{\log2}{\log(T/\pi)}\gg(\log T)^{-1}.
\]
Also $F_u'(x)\ge\log4/\log(T/\pi)\gg(\log T)^{-1}$.
For $u>4T$, the phase derivative satisfies $F_u'(x)>1$, and
\[
 F_u'(x)-1=\frac{\log(2\pi u/\tau_x)}{\log(\tau_x/2\pi)}
\ge\frac{\log(4\pi)}{\log(T/\pi)}\gg(\log T)^{-1}.
\]
Moreover $F_u'(x)\le 2b\log T/\log(T/2\pi)<2-(1-b)$ for $T$ sufficiently large. Thus in both nonresonant regions $F_u'$ stays uniformly at distance $\gg_b1/\log T$ from every integer. Kusmin--Landau bounds the unnormalized sum by $O_b(\log T)$, which gives \eqref{pp:eq:nonres}.

In the resonant region $u\in[T/(4\pi),4T]$, we have $\log u\asymp\log T$. Equation \eqref{pp:eq:derivs} implies
\[
\frac{c}{T\log^2T}\le|F_u''(x)|\le\frac{C}{T\log^2T}
\]
with absolute constants $c,C>0$, uniformly on the index interval. Apply \eqref{pp:eq:vdC} with $\lambda\asymp(T\log^2T)^{-1}$ and length $M\asymp T\log T$. The unnormalized sum is $O(\sqrt T\log T)$, giving \eqref{pp:eq:res}.
\end{proof}

\subsection{Counting resonant prime-power pairs}
Write
\begin{equation}
\mathcal R_b(T)=\#\{(n,m)\in\mathcal Q_b(T)^2:\ T/(4\pi)\le nm\le4T\}.
\label{pp:eq:R}
\end{equation}
We use only the elementary prime-count upper bound $\pi(y)\ll y/\log y$ for $y\ge3$ and the bound
\begin{equation}
\#\{p^a\le y:a\ge2\}\ll\sqrt y\qquad(y\ge2).
\label{pp:eq:powercount}
\end{equation}
Indeed the $a=2$ contribution is $\le\sqrt y$, and the remaining $a\ge3$ contribute $O(y^{1/3}\log y)=O(\sqrt y)$.

\begin{lemma}[Sparse resonant region]\label{pp:lem:pairs}
For each fixed $0<b<1$,
\begin{equation}
\mathcal R_b(T)\ll_b\frac{T}{\log T}.
\label{pp:eq:paircount}
\end{equation}
\end{lemma}
\begin{proof}
For a pair counted in \eqref{pp:eq:R}, the bounds $n,m\le T^b$ force
\begin{equation}
 n,m\ge\frac{T^{1-b}}{4\pi}.
\label{pp:eq:bothlarge}
\end{equation}
Partition the possible $n$ into $O(\log T)$ dyadic intervals $[N,2N)$, where $N\gg T^{1-b}$ and $N\ll T^b$. For each such interval, $m$ must lie in
\begin{equation}
\frac{T}{8\pi N}\le m\le\frac{4T}{N}.
\label{pp:eq:mdyad}
\end{equation}
Both $N$ and $T/N$ are bounded below by a positive power of $T$ depending on $b$. Thus the number of prime $n$ in the interval is $O_b(N/\log T)$, and the number of prime $m$ in \eqref{pp:eq:mdyad} is $O_b((T/N)/\log T)$. Consequently, prime--prime pairs contribute at most $O_b(T/\log^2T)$ on each dyadic interval, and hence
\begin{equation}
\mathcal R_{\mathrm{prime},\mathrm{prime}}\ll_b T/\log T.
\label{pp:eq:pp}
\end{equation}
If $n$ is prime and $m$ a proper prime power, the same dyadic interval contributes at most
\[
O_b\left(\frac{N}{\log T}\sqrt{\frac{T}{N}}\right)
=O_b\left(\frac{\sqrt{TN}}{\log T}\right).
\]
Summing the geometric series over dyadic $N\ll T^b$ gives
$O_b(T^{(1+b)/2}/\log T)$; reverse the roles to handle proper-power--prime pairs. If both are proper prime powers, each dyadic interval contributes at most
$O(\sqrt N\sqrt{T/N})=O(\sqrt T)$, hence $O(\sqrt T\log T)$ overall. Since $b<1$ is fixed, both bounds are $o(T/\log T)$. Combine these estimates with \eqref{pp:eq:pp}.
\end{proof}

\begin{theorem}[Prime-power almost orthogonality]\label{pp:thm:main}
For every fixed $0<b<1$, the Gram matrix in \eqref{pp:eq:gram} satisfies
\begin{equation}
\boxed{\quad\|G-I\|_{\mathrm{op}}
\ll_b T^{b-1}\log T+(\log T)^{-1/2}.\quad}
\label{pp:eq:main}
\end{equation}
In particular $\|G-I\|_{\mathrm{op}}\ll_b(\log T)^{-1/2}$, and for all sufficiently large $T$ the vectors $(v_n)_{n\in\mathcal Q_b(T)}$ are linearly independent.
\end{theorem}
The proof immediately below combines the quotient and product estimates with the resonant-pair count; it does not assume RH.

\begin{remark}[Scope]
The restriction $b<1$ is essential to the argument, both in the quotient estimates and in keeping the first derivative of product phases away from the integer $2$. There is no assertion uniform as $b\uparrow1$, nor an assertion at $b=1$. The vectors are normalized with the ordinary Euclidean scalar product, not with a second factor $1/M$.
\end{remark}



\subsection{Proof of prime-power almost orthogonality}
\begin{proof}[Proof of the prime-power almost-orthogonality theorem]
The identity $2\sin A\sin B=\cos(A-B)-\cos(A+B)$ gives
\begin{equation}
G_{nm}=\Re\,C_T(n/m)-\Re\,C_T(nm).
\label{pp:eq:decomp}
\end{equation}
Let $D_{nm}=\mathbf1_{n\ne m}\Re\,C_T(n/m)$ and $H_{nm}=\Re\,C_T(nm)$. Since $C_T(1)=1$,
\begin{equation}
G-I=D-H.
\label{pp:eq:DH}
\end{equation}
Put $N_b=|\mathcal Q_b(T)|\le T^b$. By Lemma \ref{pp:lem:quotient}, summing over all integers $m$ rather than just prime powers,
\begin{equation}
\sum_{m\in\mathcal Q_b(T):m\ne n}|D_{nm}|
\ll_b \frac{T^b}{T}\sum_{1\le |d|\le T^b}\frac1{|d|}
\ll_b T^{b-1}\log T.
\label{pp:eq:rows}
\end{equation}
The matrix $D$ is real symmetric. The Schur row-sum bound therefore gives
\begin{equation}
\|D\|_{\mathrm{op}}\ll_b T^{b-1}\log T.
\label{pp:eq:Dop}
\end{equation}

By Lemma \ref{pp:lem:product}, separating the $\mathcal R_b(T)$ resonant pairs from the others,
\begin{equation}
\|H\|_{\mathrm{HS}}^2=\sum_{n,m\in\mathcal Q_b(T)}|\Re\,C_T(nm)|^2
\ll_b \frac{N_b^2}{T^2}+\frac{\mathcal R_b(T)}{T}
\ll_b T^{2b-2}+\frac1{\log T},
\label{pp:eq:HHS}
\end{equation}
where Lemma \ref{pp:lem:pairs} supplies the final estimate. As $\|H\|_{\mathrm{op}}\le\|H\|_{\mathrm{HS}}$,
\begin{equation}
\|H\|_{\mathrm{op}}\ll_b T^{b-1}+(\log T)^{-1/2}.
\label{pp:eq:Hop}
\end{equation}
The triangle inequality applied to \eqref{pp:eq:DH}, \eqref{pp:eq:Dop} and \eqref{pp:eq:Hop} proves \eqref{pp:eq:main}. The conclusions follow since $T^{b-1}\log T=o((\log T)^{-1/2})$ for fixed $b<1$, and $\|G-I\|_{\mathrm{op}}<1$ implies positive definiteness.
\end{proof}

\subsection{True-ordinate index window}
\begin{corollary}[Change of index window]\label{pp:cor:zeros}
The almost-orthogonality theorem \ref{pp:thm:main} remains valid, with the same order of error, if $\mathcal J_T$ is replaced by the consecutive indices $\mathcal J_T^\gamma=\{k:T\le\gamma_k<2T\}$ of the positive ordinates of the nontrivial zeros of $\zeta$, counted with multiplicity, while the phases are \emph{still evaluated at the synthetic clock points} $\tau_k$.
\end{corollary}
\begin{proof}
The unconditional Riemann--von Mangoldt count $N(t)=N_0(t)+O(\log t)$ implies $k-\tfrac12-N_0(\gamma_k)=O(\log\gamma_k)$, with multiplicities included. Since $N_0'(t)\asymp\log t$ on dyadic intervals, the mean-value theorem yields $|\tau_k-\gamma_k|\ll1$ for large $k$. Consequently the two index sets differ only at the endpoints, by $O(\log T)$ indices, because $N_0$ advances by $O(\log T)$ over any bounded interval near $T$ or $2T$. Their cardinalities are both $\asymp T\log T$ and differ by $O(\log T)$. Every Gram entry changes by $O(1/T)$ (including the normalization), so the operator norm of the change is $\ll N_b/T\ll T^{b-1}$, by the row-sum bound. Absorb this in \eqref{pp:eq:main}. The later projection identity consequently applies to the changed window, with the Gram matrix and vectors normalized over that window.
\end{proof}


\subsection{The index-dependent cutoff stability theorem}
Retain the definitions of $N_0$, $\lamb$, $\tau_k$, $\gamma_k$, $\mathcal J_T$, $M$ and the normalized clock-window average from the preceding sections. The ordinates $\gamma_k$ are counted with multiplicity; no RH assumption is made.
The prime Euler phase includes \emph{all} powers of each prime below the \emph{prime} cutoff:
\begin{equation}
S_X(t)=\sum_{p\le X}\sum_{m\ge1} a_{p^m}\sin(mt\log p),
\qquad a_{p^m}=-\frac1{\pi m p^{m/2}}.
\label{ov:eq:fullphase}
\end{equation}
The construction uses \emph{no} true ordinate:
\begin{equation}
x_k=\tau_k-\frac{S_{X_k}(\tau_k)}{\lamb(\tau_k)},
\quad r_k=\lamb(\tau_k)(x_k-\gamma_k),
\quad s_k=k-\frac12-N_0(\gamma_k).
\label{ov:eq:construct}
\end{equation}
Fix $0<b<1$. Let $E_b(T)\subset\R^M$ be the linear space spanned by $\sin(\tau_k\log p^m)$ with $p^m\le T^b$ and let $P_b$ be the \emph{orthogonal projection} with respect to $M^{-1}\sum$. Define
\begin{equation}
\Omega_b[x]=\avgJ{\bigl((I-P_b)r\bigr)^2},\qquad
\Omega_b[s]=\avgJ{\bigl((I-P_b)s\bigr)^2}.
\label{ov:eq:Omega}
\end{equation}
No invertibility of the Gram matrix is needed for the orthogonal projection.

\subsection{Maximal control of a moving prime band}
\begin{lemma}[Prime maximal inequality]\label{ov:lem:maximal}
For any $2\le U<W\le T^{c_0}$, set
\begin{equation}
\mathcal Q(t)=\max_{U\le v\le W}
\left|\sum_{U<p\le v}p^{-1/2}e^{it\log p}\right|,
\qquad V=\sum_{U<p\le W}\frac1p.
\end{equation}
Then
\begin{equation}
\avgJ{\mathcal Q(\tau_k)^2}\ll_{c_0}V.
\label{ov:eq:max}
\end{equation}
The bound is independent of the number and positions of primes in the band.
\end{lemma}
\begin{proof}
Let $I=[T-1,2T+2]$. Choose a real nonnegative function
$W_T(t)=C\operatorname{sinc}^2(\eta(t-3T/2)/T)$ with sufficiently small absolute $\eta>0$ and $C$ sufficiently large. It satisfies $W_T\ge1$ on $I$, $\int_\R W_T^2\ll T$, and $\widehat W_T$ is supported in $[-2\eta/T,2\eta/T]$. Choose $\eta$ so that the support of $\widehat{W_T^2}$ is contained in $[-4\eta/T,4\eta/T]$.

Because consecutive distinct primes $p<q\le W$ satisfy $\log(q/p)\ge1/(2W)\gg T^{-c_0}$, the Fourier packets of
\[W_T(t)\sum_{U<p\le W}p^{-1/2}e^{it\log p}\]
are disjoint. Moreover every initial partial sum of these packets is a Fourier half-line truncation at a point between two consecutive packets. The Carleson--Hunt maximal theorem on $L^2(\R)$ therefore yields
\begin{equation}
\int_I\mathcal Q(t)^2dt
\ll\int_\R W_T(t)^2\left|\sum_{U<p\le W}p^{-1/2}e^{it\log p}\right|^2dt
=\left(\int_\R W_T^2\right)V\ll TV.
\label{ov:eq:QL2}
\end{equation}
The equality uses disjoint frequency support for the cross terms of $W_T^2$; equivalently, its Fourier transform vanishes at $\log(p/q)$ for $p\ne q$.
This is a literal application of the real-line Carleson--Hunt theorem:
for each permitted partial prime sum, multiplication by $W_T$ creates
Fourier packets on the real line, and their disjointness makes the
partial sum an exact Fourier half-line projection. No unproved
transference from an integer Fourier series to the frequencies
$\log p$ is being used.

The function $\mathcal Q$ is the maximum of finitely many Lipschitz moduli. Almost everywhere it satisfies
\[|\mathcal Q'(t)|\le
\max_{U\le v\le W}\left|\sum_{U<p\le v}i(\log p)p^{-1/2}e^{it\log p}\right|.
\]
Repeating the Fourier-packet argument with coefficients $(\log p)/\sqrt p$ gives
\begin{equation}
\int_I|\mathcal Q'(t)|^2dt\ll T(\log W)^2 V.
\label{ov:eq:Qprime}
\end{equation}
Successive RvM points in the window are separated by at least
$h=2\pi/\log(3T)$; the intervals $[\tau_k,\tau_k+h]$ are disjoint and contained in $I$. For every absolutely continuous $g$ on an interval $[a,a+h]$,
\begin{equation}
|g(a)|^2\le\frac2h\int_a^{a+h}|g(t)|^2dt+2h\int_a^{a+h}|g'(t)|^2dt.
\label{ov:eq:Sobolev}
\end{equation}
Sum \eqref{ov:eq:Sobolev} with $g=\mathcal Q$, use $hM\asymp T$, and insert \eqref{ov:eq:QL2}--\eqref{ov:eq:Qprime} to obtain
\[\avgJ{\mathcal Q(\tau_k)^2}\ll V\left(1+\frac{(\log W)^2}{\log^2T}\right)\ll_{c_0}V.\qedhere\]
\end{proof}

\subsection{Prime-power tails outside the projection}
\begin{lemma}[Unprojected high prime powers]\label{ov:lem:tail}
Assume $2\le X_-\le T^{c_0}$ and $0<b<1$. Then
\begin{equation}
F(t)=\sum_{p\le X_-}\sum_{m\ge2:\,p^m>T^b}a_{p^m}\sin(mt\log p)
\end{equation}
satisfies
\begin{equation}
\avgJ{F(\tau_k)^2}^{1/2}\ll_{b,c_0}T^{-b/4}\log\log(3T).
\label{ov:eq:tail}
\end{equation}
\end{lemma}
\begin{proof}
Fix $A=4$. First remove the terms $p^m>T^A$. For each $p\le T^{c_0}$, the geometric tail of $p^{-m/2}/m$, starting at the first $m$ with $p^m>T^A$, is $O(T^{-A/2})$, uniformly in $p$. Summing over at most $T^{c_0}$ primes, this contributes $O(T^{c_0-A/2})=O(T^{c_0-2})$ pointwise.

For each $m\ge2$, put
\[F_m(t)=\sum_{\substack{p\le X_-\\T^b<p^m\le T^A}}a_{p^m}\sin(mt\log p).\]
These polynomials vanish unless $m\le A\log T/\log2$, and all occurring frequencies $m\log p$ are bounded by $A\log T$. Their positive frequencies are separated by at least $m/(2X_-)$, since $m\log(q/p)\ge m/(2X_-)$ for distinct $p,q\le X_-$. The generalized Montgomery--Vaughan/Hilbert inequality for separated frequencies yields, on any interval $J$ of length $T+O(1)$,
\begin{equation}
\int_J|F_m|^2dt\ll(T+X_-/m)\sum_p|a_{p^m}|^2,
\quad
\int_J|F_m'|^2dt\ll(T+X_-/m)(A\log T)^2\sum_p|a_{p^m}|^2.
\label{ov:eq:MV}
\end{equation}
The same sampling inequality \eqref{ov:eq:Sobolev}, applied to $|F_m|$ (or directly to $F_m$), implies, as $X_-/T\ll T^{c_0-1}$,
\begin{equation}
\avgJ{F_m(\tau_k)^2}^{1/2}\ll_A \left(\sum_p|a_{p^m}|^2\right)^{1/2}.
\label{ov:eq:Fsamp}
\end{equation}
For $m=2$, $p^2>T^b$ implies
$\sum_p|a_{p^2}|^2\ll\sum_{n>T^{b/2}}n^{-2}\ll T^{-b/2}$. For $m\ge3$, use $p^{-m}\le T^{-b/2}p^{-m/2}$ on the summation range and
$\sup_{m\ge3}\sum_p p^{-m/2}<\infty$, obtaining
$\sum_p|a_{p^m}|^2\ll m^{-2}T^{-b/2}$.
Minkowski's inequality in the discrete mean and summation over $2\le m\le A\log T/\log2$ now yield
\[\avgJ{(\sum_mF_m(\tau_k))^2}^{1/2}
\ll T^{-b/4}\sum_{m\le A\log T/\log2}\frac1m
\ll T^{-b/4}\log\log(3T).\]
The discarded $O(T^{c_0-2})$ term is smaller. This proves \eqref{ov:eq:tail}.
\end{proof}

\begin{lemma}[An unconditional second moment on the clock-index window]
\label{ov:lem:ordinalmoment}
With the notation above,
\begin{equation}
 \avgJ{s_k^2}\ll\log\log T.
 \label{ov:eq:ordinalmoment}
\end{equation}
\end{lemma}
\begin{proof}
By the unconditional Riemann--von Mangoldt estimate and
\cref{lem:clocklocal}, the indices $k\in\mathcal J_T$ have
$\gamma_k\in[T-C_{\rm loc},2T+C_{\rm loc}]$ and
$|s_k|\ll\log T$; here multiplicities are included.
Apply the exact cubic counting identity \eqref{eq:cubic} with
$F=N_0$ on a slightly larger interval whose endpoints are not
ordinates. It expresses the sum of $s_k^2$ over those true ordinates
as
\[
 \int (N(t)-N_0(t))^2N_0'(t)\,dt
 +O((\log T)^3+T\log T).
\]
Selberg's unconditional continuous second moment gives
$\int_T^{2T}|S(t)|^2dt\ll T\log\log T$; the same estimate applies
to each bounded enlargement of $[T,2T]$ by splitting into a bounded
number of dyadic-length windows. Since
$N-N_0=S(t)+O(T^{-1})$ away from ordinates,
$N_0'(t)\asymp\log T$, and the endpoints contribute
$O((\log T)^3)$ using $N-N_0=O(\log T)$, the sum is
$O(T\log T\log\log T)$. Its nonnegative summands include every
$k\in\mathcal J_T$, and $M\asymp T\log T$.
This proves the claim without substituting continuous mean values
for discrete moments. Korolev's discrete Gram-point moments provide
an alternative route, but are not required here.
\end{proof}

\begin{theorem}[Arbitrary index dependence within a thin prime band]\label{ov:thm:main}
Fix $0<c_0<b<1$. For every sufficiently large $T$, choose \emph{arbitrarily}, independently at every index $k\in\mathcal J_T$, cutoffs $X_k$ satisfying
\begin{equation}
2\le X_-(T)\le X_k\le X_+(T)\le T^{c_0}.
\label{ov:eq:band}
\end{equation}
Put
\begin{equation}
V_T=\sum_{X_-<p\le X_+}\frac1p.
\label{ov:eq:V}
\end{equation}
Then, unconditionally and uniformly in the complete sequence $(X_k)$,
\begin{equation}\label{ov:eq:sqrt}
\left|\sqrt{\Omega_b[x]}-\sqrt{\Omega_b[s]}\right|
\ll_{b,c_0}\sqrt{V_T}+V_T+T^{-b/4}\log\log(3T)+T^{-1}.
\end{equation}
where $V_T=0$ is allowed. In particular, using
Lemma~\ref{ov:lem:ordinalmoment},
\begin{align}
|\Omega_b[x]-\Omega_b[s]|
&\ll_{b,c_0}\sqrt{\log\log T}\,\mathcal E_T+\mathcal E_T^2,
\\
\mathcal E_T&=\sqrt{V_T}+V_T+T^{-b/4}\log\log(3T)+T^{-1}.
\label{ov:eq:energy}
\end{align}
If $V_T\ll1/\log T$, then
\begin{equation}
\boxed{\Omega_b[x]=\Omega_b[s]
 +O_{b,c_0}\left(\sqrt{\frac{\log\log T}{\log T}}\right).}
\label{ov:eq:claimed}
\end{equation}
Crucially, $k\mapsto X_k$ need not be monotone, continuous, or representable as a smooth function of $\tau_k$.
\end{theorem}

\begin{remark}[Precisely what small harmonic mass implies]
\label{ov:rem:massrate}
As $T\to\infty$, the condition $V_T=o(1)$ implies
$|\sqrt{\Omega_b[x]}-\sqrt{\Omega_b[s]}|=o(1)$ by
\eqref{ov:eq:sqrt}.  For the energies themselves the quantitative
estimate \eqref{ov:eq:energy} proves $\Omega_b[x]-\Omega_b[s]=o(1)$
under the stronger sufficient condition
$V_T\log\log T\to0$ (assuming the displayed negligible
prime-power and Taylor tails). The three explicit cutoff laws below
satisfy the stronger bound $V_T\ll1/\log T$.
It is not correct to infer energy convergence from $V_T=o(1)$
\emph{alone} using this argument: the unprojected second moment
$\Omega_b[s]$ can grow with $T$.
\end{remark}




\begin{lemma}[Quantitative harmonic-mass criterion for arbitrary index cutoffs]
\label{ov:lem:masscriterion}
Let $2\le X_-\le X_+$, and put
\[
  u_T=\log X_-,\qquad
  \Delta_T=\log(X_+/X_-),\qquad
  H_0=\log\log T.
\]
No regularity or monotonicity of the individual $X_k$ is assumed.
Uniformly for $X_-\ge2$,
\begin{equation}
\begin{split}
 V_T
 &=\log\left(1+\frac{\Delta_T}{u_T}\right)
       +O\left(\frac1{u_T}\right),\\
 V_TH_0&\ll\frac{H_0(\Delta_T+1)}{u_T}.
\end{split}
\label{ov:eq:masscriterion}
\end{equation}
In particular, the purely deterministic condition
\begin{equation}
 \frac{H_0(\Delta_T+1)}{u_T}\longrightarrow0
 \label{ov:eq:masssuff}
\end{equation}
implies $V_T\log\log T\to0$, and therefore the energy
stability $\Omega_b[x]-\Omega_b[s]\to0$ in
Theorem~\ref{ov:thm:main}, provided its other hypotheses hold.
If $X_-\gg T^{c_-}$ and $X_+/X_-\le C$, with fixed
$c_->0$ and $C\ge1$, then quantitatively
\begin{equation}
 V_T\ll_{c_-,C}\frac1{\log T},\qquad
 V_TH_0\ll_{c_-,C}\frac{\log\log T}{\log T}.
 \label{ov:eq:massstrong}
\end{equation}
\end{lemma}
\begin{proof}
Mertens' unconditional prime-harmonic theorem gives
$\sum_{p\le y}p^{-1}=\log\log y+B+O(1/\log y)$ uniformly for
$y\ge2$. Subtract the formula at $X_-$ from that at $X_+$;
use $\log X_+=u_T+\Delta_T$ and
$0\le\log(1+z)\le z$ for $z\ge0$. Under the final hypotheses,
$u_T\ge c_-\log T+O(1)$ and $\Delta_T\le\log C$.
The energy implication is exactly \eqref{ov:eq:energy}.
\end{proof}

\begin{corollary}[Indexwise varying exponents without monotonicity]
\label{ov:cor:generalvary}
Suppose $X_k=(\tau_k/2\pi)^{c(k)}$ on $\mathcal J_T$,
where $0<c_-\le c(k)\le c_+<b<1$ are fixed bounds.
Write $d_T=\max_{k\in\mathcal J_T}c(k)
 -\min_{k\in\mathcal J_T}c(k)$. If
\begin{equation}
  d_T\log\log T\longrightarrow0,
  \label{ov:eq:osccondition}
\end{equation}
then $V_T\log\log T\to0$ and
$\Omega_b[x]-\Omega_b[s]\to0$, without any ordering assumption
on the $c(k)$. More precisely,
\begin{equation}
 V_T\log\log T
 \ll_{c_-,c_+} (d_T+1/\log T)\log\log T.
 \label{ov:eq:oscrate}
\end{equation}
\end{corollary}
\begin{proof}
Let $u_k=\log(\tau_k/2\pi)$. Across $\mathcal J_T$,
$u_k=U+O(1)$ with $U=\log(T/2\pi)$.
The difference between the largest and smallest products
$c(k)u_k$ is at most $d_TU+c_+\log 2$.
Moreover $\min_k c(k)u_k\ge c_-U$.
Insert these two bounds into \eqref{ov:eq:masscriterion}.
\end{proof}

\begin{remark}[Small harmonic mass alone does not imply the stronger rate]
\label{ov:rem:masscounter}
The condition $V_T\to0$ by itself does \emph{not} imply
$V_T\log\log T\to0$. For instance, fix $0<c<c_0<b<1$,
let $H_0=\log\log T$ and $\varepsilon_T=H_0^{-1/2}$, and
alternate indexwise between $X_k=T^c$ and
$X_k=T^{c(1+\varepsilon_T)}$ on $\mathcal J_T$.
For large $T$ both cutoffs lie below $T^{c_0}$, while
Mertens' theorem gives
\[
 V_T=\log(1+\varepsilon_T)+O_c(1/\log T)
 \sim H_0^{-1/2},\qquad V_TH_0\sim\sqrt{H_0}\to\infty.
\]
This counterexample concerns the logical implication between
conditions on the band mass; it does not assert that
$\Omega_b[x]-\Omega_b[s]$ fails to tend to zero in this example.
\end{remark}

\subsection{Proof of the stability theorem}
\begin{proof}[Proof of the index-dependent cutoff stability theorem]
Let $y_k=S_{X_k}(\tau_k)$. Split it \emph{exactly} into
\begin{equation}
\begin{split}
y_k={}&y_k^{\mathrm{proj}}+F(\tau_k)+e_k,\\
y_k^{\mathrm{proj}}={}&\sum_{\substack{p\le X_-\\p^m\le T^b}}a_{p^m}\sin(\tau_k\log p^m),\\
e_k={}&\sum_{X_-<p\le X_k}\sum_{m\ge1}a_{p^m}\sin(m\tau_k\log p).
\end{split}
\label{ov:eq:split}
\end{equation}
The vector $y^{\mathrm{proj}}$ belongs to $E_b(T)$ and is \emph{annihilated} by $I-P_b$.

By Lemma~\ref{ov:lem:maximal}, the first harmonics of $e_k$ have discrete $L^2$ norm $O(\sqrt{V_T})$. All higher harmonics in the moving band obey the \emph{pointwise} estimate
\begin{equation}
\left|\sum_{X_-<p\le X_k}\sum_{m\ge2}a_{p^m}\sin(m\tau_k\log p)\right|
\le C\sum_{X_-<p\le X_+}\frac1p=C V_T,
\end{equation}
by geometric summation of $\sum_{m\ge2}(mp^{m/2})^{-1}$. Thus
\begin{equation}
\avgJ{e_k^2}^{1/2}\ll\sqrt{V_T}+V_T.
\label{ov:eq:edge}
\end{equation}
Lemma~\ref{ov:lem:tail} controls $F$.

For $\delta_k=\tau_k-\gamma_k$, Taylor's formula for $N_0$ gives the \emph{exact} identity
\begin{equation}
\lamb(\tau_k)\delta_k=s_k+\chi_k,
\qquad \chi_k=\lamb(\tau_k)\delta_k-\bigl(N_0(\tau_k)-N_0(\gamma_k)\bigr).
\label{ov:eq:chi}
\end{equation}
The unconditional estimates $\delta_k=O(1)$ and $\lamb'(t)=1/(2\pi t)$ imply $\sup_{\mathcal J_T}|\chi_k|\ll T^{-1}$. (They follow from the classical unconditional bound $S(t)=O(\log t)$ and the Riemann--von Mangoldt count; multiplicities are permitted.) Accordingly,
\begin{equation}
r=s-y+\chi,
\qquad(I-P_b)r=(I-P_b)s-(I-P_b)(F+e-\chi).
\end{equation}
An orthogonal projection is a contraction. The reverse triangle inequality, \eqref{ov:eq:edge} and Lemma~\ref{ov:lem:tail} prove \eqref{ov:eq:sqrt}.

Finally Lemma~\ref{ov:lem:ordinalmoment} gives $\sqrt{\Omega_b[s]}\le\|s\|_2\ll\sqrt{\log\log T}$. Since
\[|\Omega_b[x]-\Omega_b[s]|\le
\left|\sqrt{\Omega_b[x]}-\sqrt{\Omega_b[s]}\right|
\left(2\sqrt{\Omega_b[s]}+
\left|\sqrt{\Omega_b[x]}-\sqrt{\Omega_b[s]}\right|\right),\]
we obtain \eqref{ov:eq:energy}--\eqref{ov:eq:claimed}.
\end{proof}

\subsection{Index-dependent exponent laws}
\begin{corollary}[Macroscopic, varying exponent]\label{ov:cor:osc}
Fix $0<a<c_*$ and $c_*+a<b<1$. Put
\begin{equation}
\boxed{c(k)=c_*+a\sin\bigl(\log\log(\tau_k/2\pi)\bigr),\qquad
X_k=(\tau_k/2\pi)^{c(k)}.}
\label{ov:eq:osc}
\end{equation}
Then $c(k)$ is nonconstant and varies through the fixed interval $[c_*-a,c_*+a]$ as $k\to\infty$. It nevertheless satisfies
\begin{equation}
\Omega_b[x]=\Omega_b[s]
+O_{a,c_*,b}\left(\sqrt{\frac{\log\log T}{\log T}}\right).
\end{equation}
\end{corollary}
\begin{proof}
Write $u=\log(t/2\pi)$ and $F(u)=u(c_*+a\sin\log u)$. Then
$F'(u)=c_*+a\sin\log u+a\cos\log u$ is bounded uniformly. The values of $u$ on $[T,2T]$ lie in an interval of length $\log2$. Consequently
$\log(X_+/X_-)\ll_{a,c_*}1$ and $X_-\ge T^{c_*-a+o(1)}$; also $X_+\le T^{c_0}$ for some $c_*+a<c_0<b$. By Mertens' prime harmonic estimate, the exact band mass satisfies
\[
 V_T=\log\frac{\log X_+}{\log X_-}
       +O\!\left(\frac1{\log X_-}\right).
\]
Here $\log(X_+/X_-)\ll_{a,c_*}1$ and
$\log X_-\asymp_{a,c_*}\log T$, whence
$V_T\ll_{a,c_*}(\log T)^{-1}$. Apply Theorem~\ref{ov:thm:main}.
\end{proof}

\begin{corollary}[Arbitrarily oscillating index cutoff]\label{ov:cor:wild}
Fix $0<c<b<1$ and $A>0$, and let $(h_k)$ be \emph{any} sequence satisfying $|h_k|\le A$ (for example $h_k=A(-1)^k$). Put
\begin{equation}
\boxed{X_k=(\tau_k/2\pi)^c e^{h_k},\qquad
c(k)=c+\frac{h_k}{\log(\tau_k/2\pi)}.}
\end{equation}
Then the conclusion \eqref{ov:eq:claimed} holds, uniformly in $(h_k)$.
\end{corollary}
\begin{proof}
Over $k\in\mathcal J_T$, $X_+/X_-\le2^ce^{2A}$ and $X_-\ge e^{-A}(T/2\pi)^c$. Choose $c<c_0<b$; for large $T$ the band satisfies \eqref{ov:eq:band}. Mertens' estimate
$V_T=\log(\log X_+/\log X_-)+O(1/\log X_-)$
therefore gives $V_T\ll_{A,c}1/\log T$, independently of
all variations of $h_k$.
\end{proof}

\begin{corollary}[First slowly varying Casa\~na cutoff]\label{ov:cor:slow}
Let
\begin{equation}
X_k=(\tau_k/2\pi)^{(\log\log\tau_k)^{-1/2}}
\end{equation}
for sufficiently large $\tau_k$. For every fixed $b\in(0,1)$, the estimate \eqref{ov:eq:claimed} holds.
\end{corollary}
\begin{proof}
The derivative of $\log X(t)$ with respect to $\log t$ is $O((\log\log t)^{-1/2})$, so $\log(X_+/X_-)\ll(\log\log T)^{-1/2}$. Also $\log X_-\asymp\log T/(\log\log T)^{1/2}$. Set $h=X_+-X_-\ll X_-/(\log\log T)^{1/2}$. For large $T$, $\log h\asymp\log X_-$. The unconditional Brun--Titchmarsh estimate $\pi(y+h)-\pi(y)\ll h/\log h$ then gives
\[
V_T\le\frac{\pi(X_+)-\pi(X_-)}{X_-}
\ll\frac{h}{X_-\log h}\ll\frac1{\log T}.
\]
Furthermore $X_+\le T^{c_0}$ for any fixed $c_0\in(0,b)$ once $T$ is large. Apply Theorem~\ref{ov:thm:main}.
\end{proof}

\begin{remark}[The actual hypothesis, and its limitation]
The conclusion is \emph{uniform for arbitrary index-by-index variation} in the band \eqref{ov:eq:band}. The smallness of the harmonic mass $V_T$, not index monotonicity, is the hypothesis. Merely assuming $X_k\le T^{c_0}$ while permitting $X_k$ to range from $2$ to $T^{c_0}$ gives $V_T\asymp\log\log T$, so this theorem does \emph{not} assert stability for that unrestricted regime. The latter would require a separate cancellation mechanism and is not established here. The statement concerns the one-step construction, not nonlinear multi-step transport.
\end{remark}

\begin{remark}[Thin bands versus exponent-scale alternation]
\label{ov:rem:contrast}
The cutoff $X_k=(\tau_k/2\pi)^c e^{h_k}$ with bounded, completely
arbitrary $h_k$ may alternate every index and nevertheless satisfies
$V_T\ll1/\log T$. No piecewise-monotonicity hypothesis is needed.
This differs sharply from allowing the cutoff to alternate between
$T^{c_1}$ and $T^{c_2}$ with fixed $0<c_1<c_2<b$:
Mertens' theorem gives for the encompassing prime band
\begin{equation}
 V_T=\sum_{T^{c_1}<p\le T^{c_2}}\frac1p
     =\log\frac{c_2}{c_1}+O_{c_1}(1/\log T).
 \label{ov:eq:wideband}
\end{equation}
The latter mass does not tend to zero. Therefore the present
stability bound does not apply. This is a genuine \emph{boundary
of the theorem}, not a proof that $\Omega_b[x]$ necessarily changes
by a prescribed amount for such an alternating law: establishing
such a limit would require controlling its correlation with
$(I-P_b)s$.
A nonlinear multi-stage correction need not remain in the same
projected phase space. Its orthogonal remainder can therefore
have construction-dependent contributions; no universal lower
bound beyond the exact projection identity is asserted for it.
\end{remark}



\subsection{An exact S2 projection identity and its arithmetic normalization}
\label{sec:S2projection}
Let $\mathcal Q_b(T)$, $G_b$, and the Euclidean-normalized columns
$v_n=(\sqrt{2/M}\sin(\tau_k\log n))_{k\in\mathcal J_T}$ be those of
\cref{pp:thm:main}. Write $V_b$ for the matrix of these columns and
\begin{equation}\label{eq:S2projection}
P_b=V_bG_b^{-1}V_b^{\mathsf T},\qquad
\mathbf r=M^{-1/2}(r_k)_{k\in\mathcal J_T},\qquad
w=V_b^{\mathsf T}\mathbf r.
\end{equation}
The projection is unchanged if the scalar product on $\R^M$ is multiplied by $M^{-1}$. For all sufficiently large $T$ the inverse exists by \cref{pp:thm:main}.

\begin{theorem}[Exact S2 decomposition and unconditional relative Gram control]\label{thm:S2exact}
For every fixed $b<1$, every real error vector $r=(r_k)_{k\in\mathcal J_T}$, and all sufficiently large $T$,
\begin{equation}\label{eq:S2exact}
\avgJ{r_k^2}=w^{\mathsf T}G_b^{-1}w+\Omega_b[r],\qquad
\Omega_b[r]:=\frac1M\|(I-P_b)r\|_{\ell^2}^2\ge0.
\end{equation}
Moreover, uniformly in $r$,
\begin{equation}\label{eq:S2relative}
\avgJ{r_k^2}
=\bigl(1+O_b((\log T)^{-1/2})\bigr)\sum_{n\in\mathcal Q_b(T)}|w_n|^2
+\Omega_b[r].
\end{equation}
This remains a relative estimate when $w$ is arbitrarily small or zero.
\end{theorem}
\begin{proof}
The columns of $V_b$ span the finite-dimensional phase subspace, so
$P_b$ is the orthogonal projection and
$\|P_b\mathbf r\|_{\ell^2}^2=w^{\mathsf T}G_b^{-1}w$.
Pythagoras gives \eqref{eq:S2exact}. By \cref{pp:thm:main},
$\|G_b-I\|_{\mathrm{op}}\le\varepsilon_T\ll_b(\log T)^{-1/2}<1/2$, and the eigenvalues of $G_b^{-1}$ lie in
$[(1+\varepsilon_T)^{-1},(1-\varepsilon_T)^{-1}]$.
Thus $w^{\mathsf T}G_b^{-1}w=(1+O_b(\varepsilon_T))\|w\|^2$, including at $w=0$.
\end{proof}

\begin{proposition}[Exact arithmetic response coefficients, including $\chi_k$]\label{prop:S2responses}
For $n=p^m$ write $a_n=-(\pi m\sqrt n)^{-1}$. Let
\begin{equation}\label{eq:S2respdef}
\begin{aligned}
\mathcal R_n^{\tau}&=\frac2{a_n}\avgJ{s_k\sin(\tau_k\log n)},\\
\mathcal A_n&=\frac2{a_n}\avgJ{S_{X_k}(\tau_k)\sin(\tau_k\log n)},\\
\mathcal C_n&=\frac2{a_n}\avgJ{\chi_k\sin(\tau_k\log n)}.
\end{aligned}
\end{equation}
For the one-step error $r_k=\lamb(\tau_k)(x_k-\gamma_k)$,
\begin{equation}\label{eq:S2w}
 w_n=\frac{a_n}{\sqrt2}(\mathcal R_n^\tau-\mathcal A_n+\mathcal C_n),
\end{equation}
and consequently
\begin{equation}\label{eq:S2arithmetic}
\avgJ{r_k^2}
=\bigl(1+O_b((\log T)^{-1/2})\bigr)
 \sum_{n\in\mathcal Q_b(T)}\frac{a_n^2}{2}
 (\mathcal R_n^\tau-\mathcal A_n+\mathcal C_n)^2
+\Omega_b[x].
\end{equation}
No hypothesis concerning a limiting prime-response law is needed.
\end{proposition}
\begin{proof}
By Taylor's theorem for $N_0$,
$\lamb(\tau_k)(\tau_k-\gamma_k)=s_k+\chi_k$, with
$\chi_k=O(T^{-1})$ on $\mathcal J_T$ by \cref{lem:clocklocal}.
Hence $r_k=s_k-S_{X_k}(\tau_k)+\chi_k$ exactly.
Multiplying by $\sqrt2\sin(\tau_k\log n)$, averaging over the clock window, and using \eqref{eq:S2respdef} gives \eqref{eq:S2w}.
Substitution in \cref{thm:S2exact} proves \eqref{eq:S2arithmetic}.
\end{proof}

\begin{corollary}[Lower bound and nonnegative irreducible component]\label{cor:S2lower}
For every one-step cutoff sequence in the range of the stability theorem,
\begin{equation}\label{eq:S2lower}
\avgJ{r_k^2}\ge\Omega_b[x]
\ge\Omega_b[s]-O_{b,c_0}(\sqrt{\log\log T}\,\mathcal E_T+\mathcal E_T^2),
\end{equation}
where $\mathcal E_T$ is defined in \eqref{ov:eq:energy}.
If $V_T\ll1/\log T$, the final error is
$O_{b,c_0}(\sqrt{\log\log T/\log T})$.
\end{corollary}
\begin{proof}
The first inequality is Pythagoras, and the second is the energy
stability bound \eqref{ov:eq:energy}.
\end{proof}

\begin{corollary}[What a conditional response law would add]\label{cor:S2conditional}
Suppose, additionally, that a proposed prime-response law supplies a
uniform summable replacement of the finite sum in
\eqref{eq:S2arithmetic}, including the moving cutoff and the
prime-power terms. Then the corresponding integral response profile
may replace that finite sum, with the error stated by that law.
Neither \cref{thm:S2exact,prop:S2responses} nor the stability of
$\Omega_b$ proves such a uniform replacement. In particular, the
asserted profile customarily called H1 is a separate assumption.
\end{corollary}
\begin{proof}
Under the stated assumption, substitute its summable approximation
into \eqref{eq:S2arithmetic}. The dependence on H1 is explicit.
\end{proof}


\subsection{Consequences for transport energy and ordinate accuracy}
\begin{proposition}[Indexwise stability of the full one-step error]\label{prop:indexfull}
Under the hypotheses of \cref{ov:thm:main}, set
$r_k^-=\lamb(\tau_k)(\tau_k-S_{X_-}(\tau_k)/\lamb(\tau_k)-\gamma_k)$.
Then
\begin{equation}\label{eq:indexfull}
\left|\sqrt{\avgJ{r_k^2}}-\sqrt{\avgJ{(r_k^-)^2}}\right|
\ll_{c_0}\sqrt{V_T}+V_T.
\end{equation}
Consequently, with $\varepsilon_T=\sqrt{V_T}+V_T$,
\begin{equation}\label{eq:indexenergy}
\left|\avgJ{r_k^2}-\avgJ{(r_k^-)^2}\right|
\ll_{c_0}\varepsilon_T\sqrt{\avgJ{(r_k^-)^2}}+\varepsilon_T^2.
\end{equation}
These bounds apply without any projection and include every choice of
cutoff at each separate index.
\end{proposition}
\begin{proof}
By the exact definitions,
$r_k-r_k^-=-(S_{X_k}(\tau_k)-S_{X_-}(\tau_k))$.
The first harmonics of this band have norm
$O(\sqrt{V_T})$ by \cref{ov:lem:maximal}; the remaining harmonics
are bounded pointwise by $O(V_T)$. Reverse triangle inequality and
$|A^2-B^2|\le|A-B|(A+B)$ complete the proof.
\end{proof}

\begin{lemma}[Boundary change between the clock and zero index windows]\label{lem:boundaryswitch}
Fix $c_0<1$. Suppose $X_k\le T^{c_0}$ on the union of the windows
$\{k:T\le\tau_k<2T\}$ and $\{k:T<\gamma_k\le2T\}$.
For one-step errors $r_k$ the normalized quadratic means on these
two windows differ by $o(1)$. The same conclusion holds for a
fixed cutoff $X\le T^{c_0}$.
\end{lemma}
\begin{proof}
By \cref{lem:clocklocal}, $|\tau_k-\gamma_k|\le C_{\rm loc}$,
so the index windows differ by $O(\log T)$ indices near their
endpoints. The midpoint count gives sizes $\asymp T\log T$ for
both windows. Also $|s_k|\ll\log T$ uniformly,
$|\chi_k|\ll T^{-1}$, and by the trivial prime-power Euler phase
bound $|S_{X_k}(\tau_k)|\ll\sqrt{X_k}\,\log(3T)$.
Consequently $\sup|r_k|^2\ll(\log T)^2(1+T^{c_0})$.
The $O(\log T)$ boundary terms divided by $T\log T$ are
$O(T^{c_0-1}\log^2T+T^{-1}\log^2T)=o(1)$.
The small normalization difference is estimated by the same
supremum bound and the relative cardinality discrepancy
$O(T^{-1})$.\end{proof}

\begin{corollary}[Transport energy for genuinely moving cutoffs]\label{cor:indexenergyfixed}
Fix $0<c\le\frac12$ and $b$ with $c<b<1$. Let
\begin{equation}\label{eq:indexcutoff}
X_k=(\tau_k/2\pi)^c\exp(h_k),\qquad |h_k|\le A,
\end{equation}
where the sequence $(h_k)$ is arbitrary and may oscillate at every
index. Then the original fixed-power transport law extends to
this fully index-dependent construction:
\begin{equation}\label{eq:indexenergyfixed}
\frac1{M_T}\sum_{k:T<\gamma_k\le2T}
[\lamb(\tau_k)(x_k-\gamma_k)]^2
=\frac{H_1}{2\pi^2}+O_{c,A}(\sqrt{H_1H_2}).
\end{equation}
In particular,
\begin{equation}\label{eq:indexdistancefixed}
\frac1{M_T}\sum_{k:T<\gamma_k\le2T}(x_k-\gamma_k)^2
=\frac{2H_1+O_{c,A}(\sqrt{H_1H_2})}{\log^2T}.
\end{equation}
\end{corollary}
\begin{proof}
On each clock window, the cutoffs lie in a multiplicative band
$[X_-,X_+]$ satisfying
\[
\frac{X_+}{X_-}\ll_{c,A}1,\qquad X_\pm\asymp_{c,A}T^c.
\]
The ordinary prime-count bound gives
$V_T\ll_{c,A}(\log T)^{-1}$.

Comparison with the common cutoff $T^c$ uses a slightly larger
band of the same type. By
\cref{prop:indexfull}, the difference of $L^2$ norms is
$O_{c,A}((\log T)^{-1/2})$.

The fixed-power theorem \cref{thm:energy} gives
$\av{r_k(T^c)^2}=H_1/(2\pi^2)+O_c(\sqrt{H_1H_2})$.

By \cref{lem:boundaryswitch} the clock- and zero-index window means
are equal up to $o(1)$. The difference of their squared norms is
$O_{c,A}(\sqrt{H_1/\log T})=o(1)$; absorb it into the displayed
error. The height-scale factor satisfies
\begin{equation*}
\lamb(\tau_k)^{-2}
=\frac{4\pi^2}{\log^2T}
\left(1+O\left(\frac1{\log T}\right)\right).
\end{equation*}
\end{proof}

\begin{corollary}[Slowly drifting or oscillating exponents]\label{cor:indexvariableleading}
Let $c(k)$ vary with $k$ and set $X_k=(\tau_k/2\pi)^{c(k)}$.
Assume that, in every clock dyadic window, the family of cutoffs
lies in $[X_-,X_+]$, where
\begin{equation}\label{eq:indexvarycond}
T^{\eta}\le X_-\le X_+\le T^{1/2+o(1)},
\qquad V_T=\sum_{X_-<p\le X_+}\frac1p\ll\frac1{\log T},
\end{equation}
for some fixed $\eta>0$, and that $X_+\le T^{c_0}$ for a fixed
$c_0<b<1$ eventually. Then
\begin{equation}\label{eq:indexvariableleading}
\avgJ{r_k^2}\sim\frac{H_1}{2\pi^2},
\end{equation}
provided $\log X_-/\log T\le1/2$ (the comparison common cutoff
is within the established analytic range).
\end{corollary}
\begin{proof}
Let $c_T=\log X_-/\log T$. The hypotheses imply
$c_T\ge\eta>0$ and $c_T\le1/2$.
By \cref{cor:varenergy}, $\av{r_k(T^{c_T})^2}\sim H_1/(2\pi^2)$;
\cref{lem:boundaryswitch} transfers the fixed-cutoff mean to
$\mathcal J_T$. By \cref{prop:indexfull}, the variable-cutoff
error has $L^2$ norm within $O((\log T)^{-1/2})$ of the common-cutoff
one, so its squared norm has the same equivalent.
\end{proof}




\section{What the spectral and counting arguments do not imply}
\label{sec:boundaries}
The exact projection and clock estimates above are unconditional. Some stronger
claims about prime responses, uniform cutoff thresholds, or zero locations do
\emph{not} follow from those estimates. We record precise obstructions as
counterexamples, not as conjectural qualifications. No real ordinate is used
in the definitions of the clocks.

\subsection{A scalar prime-band moment is not an operator-norm theorem}
\begin{proposition}[Counterexample to the scalar-to-matrix inference]
\label{prop:noscalarframe}
Knowledge of the quadratic moment of one chosen vector of prime-phase
coefficients, even with an exact diagonal main term, does not imply that the
Gram matrix of the underlying phases tends to the identity in operator norm.
In particular a scalar cross-term estimate for the full unweighted prime sum
cannot, by itself, establish almost orthogonality uniformly over all weighted
prime subsets.
\end{proposition}
\begin{proof}
For $N\ge2$ set
\[
G_N=\left(1-\frac1N\right)I_N+\frac1N\mathbf1\mathbf1^{\mathsf T}.
\]
This is a positive-definite correlation matrix with diagonal entries $1$;
its eigenvalues are $1-1/N$ with multiplicity $N-1$ and
$2-1/N$ once. Hence $\|G_N-I_N\|_{\mathrm{op}}=1-1/N\to1$.
Nevertheless the test vector $a=e_1$ satisfies
$a^{\mathsf T}G_Na=a^{\mathsf T}a=1$ exactly.
Every positive-definite correlation matrix is the Gram matrix of a family
of unit vectors. Thus exact agreement for one coefficient choice is
compatible with failure of operator-norm almost orthogonality.
\end{proof}
\begin{remark}[Two distinct mathematical results]
A resonance calculation for a particular scalar moment, such as an
asymptotic for
$\langle(S_X^{(1)}(\tau_k))^2\rangle$, does not replace the
operator-norm statement of \cref{pp:thm:main}. The latter has a
separate proof from quotient-frequency bounds, a uniform second-derivative
estimate and a Hilbert--Schmidt count of resonant \emph{prime-power} pairs.
This distinction also applies when an asymptotic scalar moment is verified
numerically.
\end{remark}

\subsection{The correct geometry of a good index}
Let $\kappa_{k-1}<\kappa_k<\kappa_{k+1}$ be an arbitrary strictly
increasing synthetic sequence, independent of the true ordinates, and put
\[
g_k^- =\kappa_k-\kappa_{k-1},\qquad
 g_k^+ =\kappa_{k+1}-\kappa_k,
\qquad m_k=\frac{\kappa_k+\kappa_{k+1}}2,
\qquad e_k=\kappa_k-\gamma_k.
\]
We call $k$ \emph{interval-good} exactly when
$\gamma_k\in(m_{k-1},m_k)$.
\begin{lemma}[Correct asymmetric good-index condition]
\label{lem:asymmetricgood}
One has the exact equivalence
\begin{equation}\label{eq:asymmetricgood}
\gamma_k\in(m_{k-1},m_k)
\quad\Longleftrightarrow\quad
-\frac{g_k^+}{2}<e_k<\frac{g_k^-}{2}.
\end{equation}
The symmetric condition
$|e_k|<\frac12\min(g_k^-,g_k^+)$ is sufficient but is
\emph{not} equivalent to \eqref{eq:asymmetricgood}.
\end{lemma}
\begin{proof}
The interval endpoints are
$m_{k-1}=\kappa_k-g_k^-/2$ and
$m_k=\kappa_k+g_k^+/2$.
Substitute $\gamma_k=\kappa_k-e_k$ to obtain
\eqref{eq:asymmetricgood}. For non-equivalence, let
$g_k^-=10$, $g_k^+=2$ and $e_k=3$. Then
$-1<3<5$ while $|3|\not<1$.
\end{proof}

\begin{theorem}[Multiplicity-preserving counting criterion]
\label{thm:goodcount}
Let $B_T$ denote the number of interval-bad indices with
$T<\gamma_k\le2T$. Count all positive ordinates with their
multiplicities. Apart from at most $O(\log T)$ boundary indices,
the number of zeros in this height window which are either off the
critical line or not simple is at most $4B_T$. In particular, writing
$M_T=N(2T)-N(T)$ and $d_T=B_T/M_T$, the fraction of zeros both on
$\Re s=1/2$ and simple is at least
\begin{equation}\label{eq:goodcount}
1-4d_T+O\left(\frac{\log T}{M_T}\right).
\end{equation}
This is a conditional counting implication: it makes no assertion that
$d_T$ is small for any given corrected clock.
\end{theorem}
\begin{proof}
Because the intervals $(m_{k-1},m_k)$ are ordered and disjoint,
two adjacent interval-good indices force
$\gamma_k<\gamma_{k+1}$. Thus every adjacent equality of ordinates
has at least one interval-bad endpoint. A given bad index is incident
to at most two adjacent equalities.

Every off-line zero has a distinct partner with the same ordinate,
by the symmetry $\beta+i\gamma\mapsto1-\beta+i\gamma$.
A multiple zero on the line also gives repetition in the ordinate
list. Hence every zero counted on the left belongs to a maximal
block of $r\ge2$ identical ordinates. Such a block has $r-1$
adjacent equalities, and $r\le2(r-1)$.
Therefore the number of zeros belonging to these blocks is at most
twice the number of equalities, hence at most four times the number
of bad indices incident to them. At either end of the height window,
the equalities may involve an index immediately outside; the
Riemann--von Mangoldt estimate gives $O(\log T)$ such indices in
a bounded-height enlargement, and the standard multiplicity bound
$O(\log T)$ covers zeros on the height boundary. This proves the
absolute count and \eqref{eq:goodcount}.
\end{proof}

\begin{proposition}[Positive mean-square error does not force bad indices]
\label{prop:floornoobstruction}
There is no implication, even for ordered lists with uniformly separated
synthetic points, from
\[
\liminf_{M\to\infty}\frac1M\sum_{k=1}^M e_k^2>0
\]
to a positive limiting proportion of interval-bad indices. Thus a
nonzero second-moment floor, on its own, does not refute a
potential $d_T\to0$ localization mechanism.
\end{proposition}
\begin{proof}
Take $\kappa_k=k$ and $\gamma_k=k+\frac14$ for every $k$.
Then $g_k^-=g_k^+=1$, $e_k=-1/4$, and every index is
interval-good. Yet $M^{-1}\sum_{k=1}^M e_k^2=1/16$ for all $M$.
This is a counterexample to the \emph{logical inference}; it is not a
claim that these sequences are zero ordinates of $\zeta$.
\end{proof}

\subsection{The precise normalization of the irreducible component}
\begin{lemma}[Projection of the exact defect]
\label{lem:exactOmegaChi}
Let $s_k^*=\lamb(\tau_k)(\tau_k-\gamma_k)=s_k+\chi_k$.
For any linear one-step correction $y\in E_b(T)$, the identity
\begin{equation}\label{eq:exactOmegaChi}
\Omega_b[r]=\Omega_b[s^*]
\end{equation}
holds exactly. The replacement by $\Omega_b[s]$ is in general
only asymptotic, with
\begin{equation}\label{eq:chiOmegabound}
\left|\sqrt{\Omega_b[s^*]}-\sqrt{\Omega_b[s]}\right|
\le\|\chi\|_{2,\mathcal J_T}\ll T^{-1}.
\end{equation}
\end{lemma}
\begin{proof}
The exact Taylor identity \eqref{ov:eq:chi} gives $r=s^*-y$.
Since $(I-P_b)y=0$, applying $I-P_b$ shows
$(I-P_b)r=(I-P_b)s^*$. The other bound follows from the
contractive property of $I-P_b$ and $\sup|\chi_k|\ll T^{-1}$.
\end{proof}

\begin{proposition}[Small band mass and absolute energy convergence]
\label{prop:massnecessitylogic}
A bound of the form
$|\sqrt{A_T}-\sqrt{B_T}|\le C\sqrt{V_T}$, together with
$B_T\ll\log\log T$, does \emph{not} imply $A_T-B_T\to0$
from the assumption $V_T\to0$ alone. The sufficient condition
$V_T\log\log T\to0$ in \cref{ov:rem:massrate} closes this
logical gap, but is not claimed necessary.
\end{proposition}
\begin{proof}
Let $H=\log\log T$, $B_T=H$, $V_T=1/H$, and
$A_T=(\sqrt H+1/\sqrt H)^2$. Then $V_T\to0$ and
$|\sqrt{A_T}-\sqrt{B_T}|=\sqrt{V_T}$, while
$A_T-B_T=2+1/H\to2$.
\end{proof}

\subsection{Endpoint phase stationary points and unproved scalar asymptotics}
The function
\[
\varphi_n(x)=\tau_x\log n-2\pi x
\]
for the smooth Riemann--von Mangoldt inverse has
$\varphi_n'(x)=0$ exactly at $\tau_x=2\pi n$.
The width in the index variable of the stationary-phase region is
of order $\sqrt n\log n$. A formula consisting of the \emph{full}
stationary Gaussian main term cannot hold with an $O(1)$ error
uniformly for all $n$ as the stationary point meets either endpoint.

\begin{lemma}[Half-Gaussian endpoint counterexample]
\label{lem:endpointfresnel}
Let $n=p^2\to\infty$ through prime squares and choose
$T=2\pi(n-1)$. Put $x_0=N_0(2\pi n)+\frac12$ and
$x_-=N_0(T)+\frac12$. Then $x_0-x_-=O(\log n)$, whereas
$|\varphi_n''(x_0)|^{-1/2}\asymp\sqrt n\log n$.
For the oscillatory integral representing the first Poisson alias,
\begin{equation}\label{eq:halfgaussian}
\int_{x_-}^{N_0(2T)+1/2}e^{i\varphi_n(x)}\,dx
= -\left(\frac12+o(1)\right)\sqrt n\log n.
\end{equation}
In particular this integral is \emph{not}
$-\sqrt n\log n+O(1)$. For every fixed $c>1/2$ the
prime square $n=p^2$ lies in the product range
$n\le X^2$ with $X=(T/2\pi)^c$ for all large $p$.
\end{lemma}
\begin{proof}
Differentiating twice, and using
$\lamb(t)=\log(t/2\pi)/(2\pi)$, gives
\[
\varphi_n''(x_0)=-\frac{2\pi}{n(\log n)^2}.
\]
The mean-value theorem applied to $N_0$ over the interval
$[2\pi(n-1),2\pi n]$ gives $x_0-x_-=O(\log n)$.
The upper endpoint is a distance $\asymp n\log n$ from $x_0$.
After the change of variables
$u=(x-x_0)\sqrt{|\varphi_n''(x_0)|}$,
the lower endpoint tends to $0$ and the upper endpoint to
$+\infty$. The usual stationary-phase theorem with a moving
endpoint therefore gives half of the complete Gaussian
contribution. Here
$\varphi_n(x_0)=2\pi n-11\pi/4$, and the Gaussian phase
$e^{-i\pi/4}$ gives the real factor $-1$. Its complete magnitude is
$\sqrt{2\pi/|\varphi_n''(x_0)|}=\sqrt n\log n$.
This proves \eqref{eq:halfgaussian}.
\end{proof}
\begin{remark}[Correct status of a product-resonance formula]
The endpoint example refutes the \emph{uniform per-product}
full-Gaussian approximation frequently used to deduce a specific
secondary cross term from pairs $pq\asymp T/(2\pi)$. It does
\emph{not} refute an averaged formula for the total cross term:
such a theorem may remain correct after a uniform Fresnel transition
and a summation over endpoint pairs. Thus an assertion of the form
$\mathrm{Cr}^{(1)}_{(T/2\pi)^c}
=\pi^{-2}\log(c/(1-c))/\log(T/2\pi)\,(1+o(1))$
for $1/2<c<1$ must be treated as unverified by the full-Gaussian
argument alone. The independent matrix theorem
\cref{pp:thm:main} needs no such scalar asymptotic.
\end{remark}

\subsection{A fixed-$b$ conditional law has no automatic endpoint limit}
\begin{proposition}[Failure of diagonal passage $b\uparrow1$]
\label{prop:bendpoint}
Suppose a family of errors $E(T,b)$ satisfies
$E(T,b)\to0$ as $T\to\infty$ for every \emph{fixed} $b<1$.
It does not follow that $E(T,b(T))\to0$ for
$b(T)\uparrow1$, nor does it follow that the projected
remainder $\Omega_b[s]$ at a fixed $b$ can be replaced by
$\lim_{u\uparrow1}\limsup_{T\to\infty}\Omega_u[s]$ with
an $o(1)$ error.
\end{proposition}
\begin{proof}
Take
\[
E(T,b)=\mathbf1_{\{(1-b)\log T\le1\}}.
\]
For each fixed $b<1$ the expression is eventually zero.
However with $b(T)=1-(2\log T)^{-1}$ it equals $1$ for
all sufficiently large $T$. The conclusion concerning the
orthogonal remainder is subject to the same non-uniform
quantifier exchange, and so needs an independent uniform bound.
\end{proof}
\begin{remark}[Conditional response and nonlinear transport]
A law for $\mathcal R_p^\tau$ on compact intervals of
$\theta_p=\log p/\log T\in(0,1)$ is not by itself a
summable weighted-$L^2$ law for the full prime response up to
$T^b$. In addition, the index-dependent cutoff response
$\mathcal A_p$ requires a uniform arithmetic estimate, including
its near-resonant products. Even if these ingredients were proved,
no assertion about the contribution of $p>T^b$, the range
$\theta\ge1$, or the nonlinear multistep remainder would follow
automatically. The proposed value $c^*\approx0.6758$
minimizes its stated \emph{conditional integral functional};
it does not prove an unconditional optimum for corrected zero ordinates.
\end{remark}

\section{Responses to moving primes}\label{sec:responses}
For any prime $p$, define the response of the bare clock by
\begin{equation}\label{eq:Rresponse}
\mathcal R_p^\tau
=\frac2{a_p}\av{s_k\sin(\tau_k\log p)},
\end{equation}
and define the explicitly computable response of the second
phase by
\begin{equation}\label{eq:Bresponse}
\mathcal B_p^{(2)}
=\frac2{a_p}\av{\bigl(Q(x_k^{(1)})-Q(\tau_k)\bigr)
\sin(\tau_k\log p)}.
\end{equation}
These are signed correlations; no assertion about individual
moving primes follows from a harmonic-band average alone.

\begin{theorem}[Harmonic cancellation of the clock response]
\label{thm:response}
Let $0<c\le1/2$ be fixed and
$\theta_T\le\theta'<c$ with $\theta'$ allowed to vary with $T$.
Put
$\mathcal J=\{p:T^{\theta'}<p\le T^c\}$ and
$\mu'=\sum_{p\in\mathcal J}p^{-1}$.
Then, uniformly in this range,
\begin{equation}\label{eq:response}
\left|\sum_{p\in\mathcal J}\frac{\mathcal R_p^\tau}{p}\right|
\ll1+\sqrt{\mu'(1+H_2)}.
\end{equation}
In particular, for $\theta'=\theta_T$ the harmonic mean of
$\mathcal R_p^\tau$ over this entire band tends to zero,
with bound $O(\sqrt{H_2/H_1})$.
\end{theorem}
\begin{proof}
By the exact finite band identity \eqref{eq:banddual}, the sum
on the left is $2\pi^2\av{sW_{\mathcal J}(\tau)}$.
Replace $s$ by $Q(x^{(1)})$ using
\cref{lem:secondreplacement} and Cauchy--Schwarz.
The cost is $O((1+\sqrt{H_2})\sqrt{\mu'})$ by
\cref{cor:massbands}.  The expansion
\cref{lem:aggregated} applies to the prime subset $\mathcal J$
uniformly, as shown in \cref{lem:aggerrors,lem:aggoffdiag}.
Its diagonal is supported only on $p\le X_2$ and is bounded
by $O(\mu'\exp[-2\pi^2\omega^2(1+o(1))])+o(1)$,
because all its primes are at least $Y=T^{\theta_T}$.
Since $\mu'\le O(H_1)$, this is $o(1)$.
Combining proves \eqref{eq:response}.
For the full band $\mu'=\log(c/\theta_T)+o(1)
=\tfrac12H_1+O(1+\log H_2)$, and division by $\mu'$
gives the stated normalized estimate.
\end{proof}

\begin{definition}[The characteristic profile of the second clock]
\label{def:profile}
Let $P=P_0+P_1$ from \eqref{eq:truncations} and let
$\mathbb E_{\rm mod}$ denote independent uniform prime phases.
For $X_1<p\le T^{1/2}$ define
\begin{equation}\label{eq:profile}
\Phi_T(p)=\av{\mathbb E_{\rm mod}
\exp\left(-i\frac{\log p}{\Lambda(\tau_k)}P\right)},
\quad\theta_p=\frac{\log p}{\log T}.
\end{equation}
The profile is a real-valued quantity computed without the actual
ordinates.
\end{definition}

\begin{lemma}[Gaussian factor with the arithmetic variance correction]
\label{lem:profile}
Uniformly for $X_1<p\le T^{1/2}$,
\begin{equation}\label{eq:profileasymp}
\Phi_T(p)
=\exp\left[-\theta_p^2
\left(H_0-\frac12H_1+O(1)\right)+O(\theta_p^4)\right]
\left(1+O\left(\frac{\theta_p^2H_0}{\log T}\right)\right).
\end{equation}
All implicit constants are absolute.
\end{lemma}
\begin{proof}
For each prime $r\le X_1$ write its contribution to
$P$ as $f_r(\Theta_r)$.  This is odd, centered, and has
$\mathbb E f_r^2=(2\pi^2r)^{-1}+O(r^{-2})$ and
$\|f_r\|_\infty\ll r^{-1/2}$.  Independence gives
$\mathbb E e^{-i\xi P}=\prod_{r\le X_1}\mathbb E\cos(\xi f_r)$.
The factors $f_r$ here are \emph{truncated} Euler phases; the
sharp maximum $\pi^{-1}\arcsin(r^{-1/2})$ of the \emph{complete}
phase is not a valid bound for an arbitrary truncated sum.
Instead the coefficientwise triangle inequality gives, uniformly in
the truncation order,
\begin{equation}\label{eq:truncatedphasebound}
 |f_r(\vartheta)|\le\frac1\pi\sum_{m\ge1}
 \frac{r^{-m/2}}m
 =-\frac1\pi\log(1-r^{-1/2})
 \le-\frac1\pi\log(1-2^{-1/2})<\frac12.
\end{equation}
Since $\log p/\Lambda(\tau_k)\le\pi+o(1)$ for
$p\le T^{1/2}$, the factors $\mathbb E\cos(\xi f_r)$
are strictly positive uniformly in $r$ and in the truncation.
In fact $|\xi f_r|\le\pi\,[-\log(1-2^{-1/2})/\pi]+o(1)
<\pi/2$, so each cosine is positive pointwise.
Their real logarithms therefore exist.
A fourth-order Taylor expansion, with the sum of fourth moments
$\sum_{r\le X_1}\mathbb E|f_r|^4\ll\sum_r r^{-2}<\infty$,
yields
\[
\log\mathbb E e^{-i\xi P}
=-\frac{\xi^2}{4\pi^2}\sum_{r\le X_1}\frac1r
 +O(\xi^2+\xi^4).
\]
The convergent prime-by-prime correction to the variance is
included in the $O(\xi^2)$ term.  Since
$\log\log X_1=H_0-\frac12H_1$ and
$\xi=\log p/\Lambda(\tau_k)
=2\pi\theta_p(1+O(1/\log T))$
uniformly in the dyadic window, the result follows after averaging
over $k$.  The variation of $\xi$ in that average has size
$O(\theta_p/\log T)$ and produces the last relative factor.
\end{proof}

\begin{lemma}[Prime powers against arbitrary prime bands]
\label{lem:discretepp}
Let $R=X_2=T^{H_1^{-1/2}}$ and let
$\mathcal J\subset(X_1,T^{1/2}]$ be any set of primes.
Writing
\begin{equation}\label{eq:pptest}
 Q_{\mathrm{pp}}^\flat(t)
=\sum_{p\le R}\sum_{\substack{m\ge2\\p^m\le R^2}}
\beta_{p^m}\sin(t\log p^m),\qquad
 W_{\mathcal J}(t)=\sum_{q\in\mathcal J}a_q\sin(t\log q),
\end{equation}
we have, uniformly in $\mathcal J$,
\begin{equation}\label{eq:ppcorrelation}
\bigl|\langle Q_{\mathrm{pp}}^\flat(\tau)
W_{\mathcal J}(\tau)\rangle_T\bigr|
\ll T^{-1/4+o(1)}=o(1).
\end{equation}
For the complete $Q=S_R$,
\begin{equation}\label{eq:ppfullcorr}
\bigl|\langle Q_{\mathrm{pp}}(\tau)
W_{\mathcal J}(\tau)\rangle_T\bigr|
\ll T^{-1/4+o(1)}+
\frac{\sqrt{\mu(\mathcal J)}}{\log R}
=o(1+\sqrt{\mu(\mathcal J)}).
\end{equation}
\end{lemma}
\begin{proof}
Write $C_T(v)=M_T^{-1}\sum_{k\in\mathcal I_T}e^{i\tau_k\log v}$.
For positive integers $a,b\le T^{1/2+o(1)}$ with
$\max(a,b)\le T^{1/2}$ and $a\ne b$, the clock
estimate in \cref{lem:klclock}, applied to
$\log(a/b)$, yields
\begin{equation}\label{eq:clockratioPP}
|C_T(a/b)|\ll\frac{\max(a,b)}{T|a-b|}+\frac1{T\log T}
\ll\frac{\max(a,b)}T.
\end{equation}
Here $a=p^m\le R^2=T^{o(1)}$ and
$b=q\le T^{1/2}$; for all sufficiently large $T$,
$|\log(a/b)|\le(\tfrac12+o(1))\log T<\tfrac34\log T$.
Moreover $a\ne b$, since $m\ge2$ and $b$ is prime.
For the product frequencies $ab=p^mq\le T^{1/2+o(1)}$,
the same clock lemma gives
$|C_T(ab)|\ll T^{-1}$: their frequency is nonzero
and at most $\tfrac34\log T$, with
$\log(ab)\ge\log4$.
Expanding $2\sin(\tau\log a)\sin(\tau\log b)$
into the ratio and product cosines, and using the
coefficient bounds, we obtain
\begin{align*}
&|\langle Q_{\mathrm{pp}}^\flat(\tau)
 W_{\mathcal J}(\tau)\rangle_T|\\
&\quad\ll\frac1T
\sum_{p\le R}\sum_{\substack{m\ge2\\p^m\le R^2}}
\frac1{m p^{m/2}}
\sum_{q\in\mathcal J}\frac{p^m+q+1}{\sqrt q}\\
&\quad\ll\frac1T\left[
\biggl(\sum_{q\le T^{1/2}}q^{-1/2}\biggr)
\biggl(\sum_{p,m\ge2:p^m\le R^2}\frac{p^{m/2}}m\biggr)
+\biggl(\sum_{q\le T^{1/2}}q^{1/2}\biggr)
\biggl(\sum_{p,m\ge2:p^m\le R^2}\frac1{mp^{m/2}}\biggr)
\right].
\end{align*}
Elementary comparisons with all integers give
$\sum_{q\le T^{1/2}}q^{-1/2}\ll T^{1/4}$
and $\sum_{q\le T^{1/2}}q^{1/2}\ll T^{3/4}$.
The first prime-power sum is $R^{O(1)}=T^{o(1)}$,
and the second is $O(1+\log\log R)$:
for $m=2$ it is at most $\frac12\sum_{p\le R}p^{-1}$,
while the higher harmonics have bounded total sum.
It follows that the expression is $T^{-1/4+o(1)}$.
The bound is independent of the choice of $\mathcal J$,
including choices made after inspecting the response signs.

Finally $Q_{\mathrm{pp}}-Q_{\mathrm{pp}}^\flat=Q-Q^\flat$,
with uniform norm $O(1/\log R)$ by \eqref{eq:derivativetail}.
By the prime-band almost orthogonality in
\cref{cor:massbands},
$\|W_{\mathcal J}(\tau)\|_{2,T}\ll\sqrt{\mu(\mathcal J)}$.
Cauchy--Schwarz gives \eqref{eq:ppfullcorr}. No
continuous-time mean-value theorem is substituted for
a discrete clock estimate in this argument.
\end{proof}

\begin{theorem}[Second-clock response on arbitrary prime subsets]
\label{thm:secondresponse}
For every possibly $T$-dependent set of primes
$\mathcal J\subset(X_1,T^{1/2}]$, write
$\mu(\mathcal J)=\sum_{p\in\mathcal J}p^{-1}$.  Uniformly in
\emph{all} such sets,
\begin{equation}\label{eq:secondresponse}
\sum_{p\in\mathcal J}\frac{\mathcal B_p^{(2)}}p
=\sum_{\substack{p\in\mathcal J\\p\le X_2}}
\frac{\Phi_T(p)-1}{p}
 +o\bigl(1+\sqrt{\mu(\mathcal J)}\bigr).
\end{equation}
In particular the signed harmonic response of the second step in
subsets of $(X_2,T^{1/2}]$ has size
$o(1+\sqrt{\mu(\mathcal J)})$.
\end{theorem}
\begin{proof}
By \eqref{eq:bandb}, the left-hand side equals
$2\pi^2\av{[Q(x^{(1)})-Q(\tau)]W_{\mathcal J}(\tau)}$.
For the first term, repeat
\cref{lem:finitecorr,lem:aggregated,lem:aggerrors,lem:aggoffdiag}
with $W_{\mathcal J}$ instead of $W_{\rm pr}$.
Its diagonal is
\[
\frac1{2\pi^2}\sum_{\substack{p\in\mathcal J\\p\le X_2}}
\frac{\Phi_T(p)}p
 +o(1+\sqrt{\mu(\mathcal J)}).
\]
The truncated characteristic averages and the full $\Phi_T$
differ by at most the factorial remainder estimates of
\cref{lem:aggerrors}, uniformly for all $p$.
The total variation of the error across the entire prime range
is negligible, since $\mu(\mathcal J)\ll H_0$ and the
factorial remainders decay faster than any power of $H_0$.

For the second term $Q(\tau)$, unique factorization and
\cref{lem:orthprimes} give the prime diagonal
$(2\pi^2)^{-1}\sum_{p\in\mathcal J,p\le X_2}p^{-1}$,
with an error $o(1+\sqrt{\mu(\mathcal J)})$.
The higher prime harmonics in $Q$ have no diagonal
against prime tests. Their \emph{discrete} clock correlations
are $o(1+\sqrt{\mu(\mathcal J)})$ by
\cref{lem:discretepp}, including the omitted tail and
uniformity over all test subsets. The estimate explicitly
applies the clock exponential-sum bound to
$\log(p^m/q)$ and $\log(p^mq)$; it does not
replace discrete means by Montgomery--Vaughan integrals.
Subtract the two formulas and multiply by $2\pi^2$.
Every bound was obtained using the mass and support of the test
polynomial, not the contiguity of its primes.  Thus the error
is uniform over subsets that may depend on the signs of their
response coefficients.
\end{proof}

\begin{corollary}[A harmonic total-variation implication]
\label{cor:variation}
Suppose a collection $\mathscr P_T\subset(X_1,T^{1/2}]$
has $\mu(\mathscr P_T)\le C$ for fixed $C$.
Then the errors in \eqref{eq:secondresponse} satisfy
\begin{equation}\label{eq:variation}
\sum_{p\in\mathscr P_T}\frac{1}{p}
\left|\mathcal B_p^{(2)}
 -\mathbf1_{p\le X_2}(\Phi_T(p)-1)\right|=o(1).
\end{equation}
\end{corollary}
\begin{proof}
Apply the same uniform subset estimate to the set of positive
errors and to the set of negative errors.  They are admissible
subsets of $\mathscr P_T$, even though the choice depends on
$T$ and on the responses.  Since both harmonic masses are
$O(1)$, each signed sum is $o(1)$; their difference equals the
sum of absolute errors.
\end{proof}



\section{Nonlinear spectral transport}\label{sec:spectral}
Let $a_k,b_k$ be arbitrary finite families of real points and fix $u\in\R$. The correct nonlinear quantity is
\begin{equation}\label{eq:nonlinear}
\mathcal V(u)=\sum_k\e^{iua_k}\frac{\e^{iu(b_k-a_k)}-1}{iu},
\qquad \mathcal V(0)=\sum_k(b_k-a_k).
\end{equation}
\begin{proposition}[Exact phase identity and quadratic error]\label{prop:nonlinear}
For every finite family,
\begin{equation}\label{eq:nonlinearexact}
\mathcal V(u)=\sum_k\int_{a_k}^{b_k}\e^{iut}\,dt,
\qquad
\left|\mathcal V(u)-\sum_k(b_k-a_k)\e^{iua_k}\right|
\le\frac{|u|}{2}\sum_k|b_k-a_k|^2.
\end{equation}
In particular, with $a_k=\gamma_k$ and $b_k=\tau_k^X$, all the terms in the formula have exact finite-height meaning, irrespective of collisions or ordering.
\end{proposition}
\begin{proof}
The integral identity is direct for $u\ne0$, and follows at $u=0$ by continuity. Write the difference as
$\sum_k\e^{iua_k}\int_0^{b_k-a_k}(\e^{iuv}-1)\,dv$. The elementary estimate $|\e^{iv}-1|\le|v|$ bounds each term by $|u||b_k-a_k|^2/2$. Add these bounds.
\end{proof}

\begin{lemma}[Exponential-integral transport]\label{lem:Ei}
Let $u>0$ and $a,b\in\R$. Along the vertical segment from $\frac12+iau$ to $\frac12+ibu$, the principal exponential integral is analytic, and
\begin{equation}\label{eq:Ei}
\operatorname{Ei}((\tfrac12+ib)u)-\operatorname{Ei}((\tfrac12+ia)u)
=i\int_a^b\frac{\e^{(1/2+it)u}}{1/2+it}\,dt.
\end{equation}
Consequently, on a dyadic band, the difference of the two truncated spectral sums is a finite signed oscillatory integral without any exchange of infinite limits.
\end{lemma}
\begin{proof}
The segment lies in the open right half-plane and avoids the branch cut. Differentiate $\operatorname{Ei}((1/2+it)u)$ using $\operatorname{Ei}'(z)=\e^z/z$ and the chain rule. Its derivative in $t$ is $i\e^{(1/2+it)u}/(1/2+it)$. Integrate from $a$ to $b$. Finite summation is legitimate without conditions.
\end{proof}

\begin{remark}[Infinite reconstruction]
Passing from \cref{lem:Ei} to an infinite explicit formula involves truncation, the normalization at ordinates and the summability of signed dyadic contributions. Neither a bounded second moment of point displacements nor the local inequality in \cref{prop:nonlinear} alone justifies exchanging $T\to\infty$, $K\to\infty$ and $x\to\infty$. The required global tail estimate is a separate analytic assertion.
\end{remark}


\section{Separation and spectral rigidity}\label{sec:separation}
\begin{theorem}[A direct adjacent-gap certificate]\label{thm:gap}
Let $c_1<c_2<\cdots$ be a strictly increasing real sequence and let $\gamma_1\le\gamma_2\le\cdots$ be the full list of positive ordinates of the nontrivial zeros of $\zeta$, counted with multiplicity. If, for all $k$,
\begin{equation}\label{eq:gap}
|c_k-\gamma_k|+|c_{k+1}-\gamma_{k+1}|<c_{k+1}-c_k,
\end{equation}
then all positive ordinates are distinct. Consequently every nontrivial zero is simple and has real part $1/2$.
\end{theorem}
\begin{proof}
The triangle inequality gives
$\gamma_{k+1}-\gamma_k\ge c_{k+1}-c_k-|\gamma_{k+1}-c_{k+1}|-|\gamma_k-c_k|>0$.
If a nontrivial zero $\beta+i\gamma$ had $\beta\ne1/2$, then $1-\beta+i\gamma$ would also be a zero by the functional equation and conjugation symmetry; both appear as distinct entries with the same positive ordinate. A repeated zero on the critical line likewise supplies repeated entries. Both alternatives contradict strict increase.
\end{proof}

\begin{remark}
A mean-square bound, even one tending to zero in units of the mean spacing, cannot by itself establish \eqref{eq:gap} for every $k$: exceptional indices require a separate uniform argument. Furthermore sorting a proposed corrected sequence proves only that its sorted values are ordered; it does not prove that the pointwise correction map is monotone or that sorted index labels match the actual ordinates.
\end{remark}


\section{Conclusions}\label{sec:conclusions}
The construction $\tau_k^X=\tau_k-S_X(\tau_k)/\Lambda(\tau_k)$
is determined by the smooth ordinal clock and primes.  The exact
identity for its normalized error separates the residual prime
information at the actual ordinate from the nonlinear transport of
the prime phase.  This separation is valid for all nontrivial
zeta zeros counted with multiplicity and does not assume that their
real parts equal one half.

The first quantitative conclusion is the uniform ordinal defect law
\eqref{eq:primedefect}, with endpoint $X=T^{1/2}$ included.  The
second is the fixed-exponent transport-energy formula
\eqref{eq:energy}: the transport energy grows like
$(2\pi^2)^{-1}\log\log\log T$, rather than remaining bounded,
for every fixed $0<c\le1/2$.  The mean-square error in
ordinate height has the expansion \eqref{eq:heightdistance}.
The uniform moving-cutoff equivalent \eqref{eq:varenergy} holds
under its explicit logarithmic restriction; the intermediate
lower bound \eqref{eq:lowerenergy} has a distinct, wider range.

The prime-response laws \eqref{eq:response} and
\eqref{eq:secondresponse} quantify harmonic cancellation on
moving bands.  The latter remains uniform for arbitrary
prime subsets, so bounded-mass subsets admit the harmonic
absolute-value consequence \eqref{eq:variation}.
The calculated characteristic profile retains the
$-\frac12\log\log\log T$ variance correction in
\eqref{eq:profileasymp}; it cannot be replaced uniformly
by a Gaussian factor having only $\log\log T$ in the exponent.

The finite-band spectral identity \eqref{eq:nonlinearexact}
identifies the additional signed cancellation needed to compare
prime-counting or Mertens transforms from distinct ordinate systems.
A universal adjacent-gap estimate would enforce distinct ordinates
and, by the functional-equation symmetry of the full zeta zero set,
both critical-line location and simplicity.  Neither a second-moment
estimate nor a synthetic zeta with prescribed critical-line zeros
alone provides that universal adjacent-gap estimate.


The operator-norm projection theorem \cref{pp:thm:main} supplies an
unconditional prime-power frame below every fixed exponent $b<1$.
The exact S2 identity \eqref{eq:S2arithmetic} includes the quadratic
normalization correction $\chi_k$ and retains the nonnegative residual
$\Omega_b[x]$. For index-dependent cutoffs constrained to a band satisfying
$V_T\log\log T\to0$, \cref{ov:thm:main} proves
$\Omega_b[x]=\Omega_b[s]+o(1)$, without monotonicity of the cutoff
or corrected points. Furthermore \cref{cor:indexenergyfixed} transports
the fixed-power mean-square law to arbitrary bounded index-to-index
oscillations of the cutoff, with the same leading term and error order.
These geometric and quadratic results do not establish the separate
response hypothesis H1 or a global Mertens estimate. The explicit counterexamples
and the logical boundaries of the projection method are proved in
\cref{sec:boundaries}.

\begin{thebibliography}{99}
\bibitem{Carleson1966} L. Carleson, On convergence and growth of partial sums of Fourier series, \emph{Acta Math.} \textbf{116} (1966), 135--157.
\bibitem{Hunt1968} R. Hunt, On the convergence of Fourier series, \emph{Orthogonal Expansions and Their Continuous Analogues} (1968), 235--255.
\bibitem{GrahamKolesnik1991}
S. W. Graham and G. Kolesnik, \emph{Van der Corput's Method of Exponential Sums},
London Mathematical Society Lecture Note Series 126, Cambridge University Press, 1991,
doi:10.1017/CBO9780511661976.
\bibitem{Hoeffding1963}
W. Hoeffding, Probability inequalities for sums of bounded random variables,
\emph{J. Amer. Statist. Assoc.} \textbf{58} (1963), 13--30,
doi:10.1080/01621459.1963.10500830.
\bibitem{Mertens1874}
F. Mertens, Ein Beitrag zur analytischen Zahlentheorie,
\emph{J. Reine Angew. Math.} \textbf{78} (1874), 46--62,
doi:10.1515/crll.1874.78.46.
\bibitem{Selberg1946}
A. Selberg, Contributions to the theory of the Riemann zeta-function,
\emph{Arch. Math. Naturvid.} \textbf{48} (1946), 89--155.
\bibitem{Kinnunen1997}
J. Kinnunen, The Hardy--Littlewood maximal function of a Sobolev function,
\emph{Israel J. Math.} \textbf{100} (1997), 117--124.
\bibitem{Stein1970}
E. M. Stein, \emph{Singular Integrals and Differentiability Properties of Functions},
Princeton University Press, 1970.
\bibitem{Fujii1974} A. Fujii, On the zeros of Dirichlet $L$-functions. I, \emph{Trans. Amer. Math. Soc.} \textbf{196} (1974), 225--235, doi:10.1090/S0002-9947-1974-0349603-2.
\bibitem{Fujii1981} A. Fujii, On the zeros of Dirichlet $L$-functions. II (with corrections to \emph{On the zeros of Dirichlet $L$-functions. I} and the subsequent papers), \emph{Trans. Amer. Math. Soc.} \textbf{267} (1981), 33--40.
\bibitem{MV1974} H. L. Montgomery and R. C. Vaughan, Hilbert's inequality, \emph{J. London Math. Soc.} (2) \textbf{8} (1974), 73--82, doi:10.1112/jlms/s2-8.1.73.
\bibitem{Korolev2012} M. A. Korolev, Gram's law and the argument of the Riemann zeta function, \emph{Publ. Inst. Math. (Beograd)} \textbf{92}(106) (2012), 53--78, doi:10.2298/PIM1206053K.
\bibitem{Gonek2012} S. M. Gonek, Finite Euler products and the Riemann hypothesis, \emph{Trans. Amer. Math. Soc.} \textbf{364} (2012), 2157--2191, doi:10.1090/S0002-9947-2011-05546-7.
\bibitem{Titchmarsh1986} E. C. Titchmarsh, \emph{The Theory of the Riemann Zeta-Function}, 2nd ed., revised by D. R. Heath-Brown, Oxford University Press, 1986.
\end{thebibliography}
\end{document}
